\documentclass[10pt, letterpaper]{article}

% Inhaltsverzeichnis für Pakettypen (nur für Übersicht im Header, wird nicht im Dokument angezeigt)
% 1. Seitenlayout und Ränder
% 2. Sprache und Zeichensatz
% 3. Mathematik und Theorem-Umgebungen
% 4. Eigene Makros
% 5. Diagramme und Grafiken
% 6. Tabellen und Aufzählungen
% 7. Inhaltsverzeichnis
% 8. Abschnittsüberschriften
% 9. Abstrakt-Umgebung
% 10. Todos/Notizen
% 11. Rahmen/Box-Umgebungen
% 12. Python-Integration
% 13. Literaturverwaltung
% 14. Hyperlinks
% 15. Absatzeinstellungen
% 16. Umgebungen
% 17  Graphik
% 18  Extra
% 00. Titel und Autor

\usepackage{slashed}

% --- 1. Seitenlayout und Ränder ---
\usepackage[margin=3cm]{geometry}

% --- 2. Sprache und Zeichensatz ---
\usepackage[english]{babel}
\usepackage[T1]{fontenc}
\usepackage[utf8]{inputenc}

% --- 3. Mathematik und Theorem-Umgebungen ---
\usepackage{amsmath, amssymb, amsthm}
\usepackage{mathrsfs}
\DeclareMathOperator{\WF}{WF}

% --- 4. Eigene Makros ---
\usepackage{xcolor}
\newcommand{\SKP}{\langle\cdot,\cdot\rangle}
\newcommand{\id}{\operatorname{id}}
\newcommand{\sgn}{\operatorname{sgn}}
\newcommand{\Ad}{\operatorname{Ad}}
\newcommand{\ad}{\operatorname{ad}}
\newcommand{\R}{\mathbb{R}}
\newcommand{\N}{\mathbb{N}}
\newcommand{\Q}{\mathbb{Q}}
\newcommand{\Z}{\mathbb{Z}}
\newcommand{\C}{\mathbb{C}}
\newcommand{\QU}{\mathbb{H}}
\newcommand{\Cinf}{C^{\infty}}
\newcommand{\Mod}{\mathrm{Mod}}
\newcommand{\Alg}{\mathrm{Alg}}
\newcommand{\Aut}{\mathrm{Aut}}
\newcommand{\End}{\mathrm{End}}
\newcommand{\Hom}{\mathrm{Hom}}
\newcommand{\Cl}{\mathrm{Cl}}
\newcommand{\entwurf}[1]{\textcolor{red}{#1}}
\newcommand{\GL}{\mathrm{GL}}
\newcommand{\Mat}{\mathrm{Mat}}
\newcommand{\SL}{\mathrm{SL}}
\newcommand{\SO}{\mathrm{SO}}
\newcommand{\SU}{\mathrm{SU}}
\newcommand{\Sp}{\mathrm{Sp}}
\newcommand{\Spin}{\mathrm{Spin}}
\newcommand{\Pin}{\mathrm{Pin}}
\newcommand{\Lor}{\mathrm{Lor}}
\newcommand{\Ogroup}{\mathrm{O}}
\newcommand{\U}{\mathrm{U}}

% --- 5. Diagramme und Grafiken ---
\usepackage{graphicx}
\graphicspath{{../../Library/Images/}}
\usepackage[export]{adjustbox}
\usepackage{tikz}
\usetikzlibrary{decorations.pathreplacing, arrows.meta, positioning}
\usepackage{tikz-cd}

% --- 6. Tabellen und Aufzählungen ---
\usepackage{enumitem}
\setlist[itemize]{left=0.5cm}

\newenvironment{romanenum}[1][]
  {%
    \ifx&#1&
    \else
      \textbf{#1}\quad
    \fi
    \begin{enumerate}[label=\roman*)]
  }
  {%
    \end{enumerate}%
  }

% --- 7. Inhaltsverzeichnis ---
\usepackage{tocloft}
\renewcommand{\cftsecfont}{\footnotesize}
\renewcommand{\cftsubsecfont}{\footnotesize}
\renewcommand{\cftsubsubsecfont}{\footnotesize}
\renewcommand{\cftsecpagefont}{\footnotesize}
\renewcommand{\cftsubsecpagefont}{\footnotesize}
\renewcommand{\cftsubsubsecpagefont}{\footnotesize}
\usepackage{etoc}

% --- 8. Abschnittsüberschriften ---
\usepackage{titlesec}
\titleformat{\section}{\normalfont\large\bfseries}{\thesection}{1em}{}
\titleformat{\subsection}{\normalfont\normalsize\bfseries}{\thesubsection}{0.5em}{}
\titleformat{\subsubsection}{\normalfont\normalsize\bfseries}{\thesubsubsection}{0.5em}{}
\setcounter{secnumdepth}{4}

% --- 9. Abstrakt-Umgebung ---
\usepackage{changepage}
\renewenvironment{abstract}
  {
    \begin{adjustwidth}{1.5cm}{1.5cm}
    \small
    \textsc{Abstract. –}%
  }
  {
    \end{adjustwidth}
  }

% --- 10. Todos/Notizen ---
\usepackage{todonotes}
% Hellblau definieren
\definecolor{hellblau}{rgb}{0.3,0.6,1}
\newenvironment{note}{\par\begingroup\color{hellblau}\itshape}{\par\endgroup}

% --- 11. Rahmen/Box-Umgebungen ---
\usepackage{mdframed}
\usepackage{tcolorbox}
\colorlet{shadecolor}{gray!25}

\newenvironment{customTheorem}
  {\vspace{10pt}%
   \begin{mdframed}[
     backgroundcolor=gray!20,
     linewidth=0pt,
     innertopmargin=10pt,
     innerbottommargin=10pt,
     skipabove=\dimexpr\topsep+\ht\strutbox\relax,
     skipbelow=\topsep,
   ]}
  {\end{mdframed}
   \vspace{10pt}%
  }

% --- 12. Python-Integration ---
% (Deaktiviert in dieser Version, aktiviere bei Bedarf)
% \usepackage{pythontex}
% \usepackage[makestderr]{pythontex}

% --- 13. Literaturverwaltung ---
\usepackage{csquotes}
\usepackage[backend=biber, style=alphabetic, citestyle=alphabetic]{biblatex}
\addbibresource{bibliography.bib}

% --- 14. Hyperlinks ---
\usepackage{hyperref}
\hypersetup{
  colorlinks   = true,
  urlcolor     = blue,
  linkcolor    = blue,
  citecolor    = blue,
  frenchlinks  = true
}

% --- 15. Absatzeinstellungen ---
\usepackage[parfill]{parskip}
\sloppy

% --- 16. Umgebungen ---
\usepackage{thmtools}

\newcommand{\CustomHeading}[3]{%
  \par\medskip\noindent%
  \textbf{#1 #2} \textnormal{(#3)}.\enskip%
}

\newenvironment{DEF}[2]{\begin{unitbox}\CustomHeading{Definition}{#1}{#2}}{\end{unitbox}}
\newenvironment{PROP}[2]{\begin{unitbox}\CustomHeading{Proposition}{#1}{#2}}{\end{unitbox}}
\newenvironment{THEO}[2]{\begin{unitbox}\CustomHeading{Theorem}{#1}{#2}}{\end{unitbox}}
\newenvironment{LEM}[2]{\begin{unitbox}\CustomHeading{Lemma}{#1}{#2}}{\end{unitbox}}
\newenvironment{KORO}[2]{\begin{unitbox}\CustomHeading{Corollar}{#1}{#2}}{\end{unitbox}}
\newenvironment{REM}[2]{\begin{unitbox}\CustomHeading{Remark}{#1}{#2}}{\end{unitbox}}
\newenvironment{EXA}[2]{\begin{unitbox}\CustomHeading{Example}{#1}{#2}}{\end{unitbox}}
\newenvironment{STUD}[2]{\begin{unitbox}\CustomHeading{Study}{#1}{#2}}{\end{unitbox}}
\newenvironment{CONC}[2]{\begin{unitbox}\CustomHeading{Concept}{#1}{#2}}{\end{unitbox}}
\newenvironment{OTH}[2]{\begin{unitbox}\CustomHeading{Other}{#1}{#2}}{\end{unitbox}}
\newenvironment{EXE}[2]{\begin{unitbox}\CustomHeading{Exercise}{#1}{#2}}{\end{unitbox}}
\newenvironment{MOT}[2]{\begin{unitbox}\CustomHeading{Motivation}{#1}{#2}}{\end{unitbox}}
\newenvironment{PROOF}[2]{\begin{unitbox}\CustomHeading{Proof}{#1}{#2}}{\end{unitbox}}

% --- Unit Umgebung für Source-Inhalte ---
\usepackage{mdframed}
\newmdenv[
  linewidth=1pt,
  topline=false,
  bottomline=false,
  rightline=false,
  leftmargin=0cm,
  rightmargin=0cm,
  skipabove=10pt,
  skipbelow=10pt,
  innertopmargin=0.5\baselineskip,
  innerbottommargin=0.5\baselineskip,
  backgroundcolor=gray!10,
  linecolor=gray
]{unitbox}

\newenvironment{unit}[1]
  {\begin{unitbox}\textbf{Unit #1}\par\smallskip}
  {\end{unitbox}}

% --- 17. ... ---

% --- 18. Extras ---
\usepackage{stmaryrd}
\usepackage{bbold}  % falls du \mathbb{1} nutzen willst

% --- 00. Titel und Autor ---
\title{Mein Titel}
\author{Tim Jaschik}
\date{\today}

\begin{document}

\maketitle
\rule{\textwidth}{0.5pt}
\begin{abstract}
Kurze Beschreibung …
\end{abstract}
\rule{\textwidth}{0.5pt}
\vspace{0.5cm}

\tableofcontents

\pagebreak

\subsection{Formelle Ausarbeitung des abstrakten Modellrahmens (I)}

\emph{Ziel.}
Wir präzisieren die Grundkomponenten eines allgemeinen physikalischen Modells
in einer rein strukturellen Formulierung. Die Entwicklung orientiert sich an den
Punkten (1)--(5) des vorangegangenen Schemas und stellt alle Objekte und
Morphismen explizit als Kategorien oder Funktoren dar.

\bigskip
\emph{1. Arenen und äußere Symmetrien.}

\textbf{(a) Kategorie der Arenen.}
Wir definieren eine Unterkategorie
\[
\mathbf{Geom} \subset \mathbf{Man},
\]
deren Objekte glatte Mannigfaltigkeiten $M$ mit zusätzlicher geometrischer
Struktur sind (z.\,B.\ Metrik, Orientierung, Spin-Struktur, Volumenform),
und deren Morphismen Struktur-erhaltende glatte Abbildungen sind.

Ein typischer morphismus-spezifischer Begriff:
\[
\mathrm{Hom}_{\mathbf{Geom}}(M_1,M_2)
:=\{f:M_1\to M_2\mid f^\ast(g_2)=g_1,\text{ ggf.\ weitere Strukturerhaltung}\}.
\]
Damit ist $\mathbf{Geom}$ eine (ggf.\ große) Unterkategorie von $\mathbf{Man}$,
die die kinematisch zulässigen Arenen und Koordinatenwechsel enthält.

\textbf{(b) Externe Symmetriegruppen.}
Sei $G_{\mathrm{ext}}$ eine Lie-Gruppe mit einer glatten Linksaktion auf $\mathcal A\in\mathbf{Geom}$:
\[
\alpha_{\mathrm{ext}}: G_{\mathrm{ext}}\times \mathcal A \to \mathcal A,
\quad (\Lambda,x)\mapsto \Lambda\cdot x.
\]
Diese Aktion definiert einen Funktor
\[
\alpha_{\mathrm{ext}}: BG_{\mathrm{ext}} \to \mathbf{Geom},\qquad
(\bullet\mapsto \mathcal A,\; \Lambda\mapsto\alpha_{\mathrm{ext}}(\Lambda)).
\]
Wir identifizieren $\alpha_{\mathrm{ext}}$ mit einem Gruppenhomomorphismus
\[
G_{\mathrm{ext}}\longrightarrow \Aut_{\mathbf{Geom}}(\mathcal A),
\]
wobei $\Aut_{\mathbf{Geom}}(\mathcal A)$ die Gruppe der invertierbaren Morphismen
$\mathcal A\to \mathcal A$ in $\mathbf{Geom}$ ist.

\textbf{(c) Hintergrundkategorie.}
Aus der Wirkung von $G_{\mathrm{ext}}$ erzeugt man eine kleine Unterkategorie
\[
\mathbf{Back}(\mathcal A):
=\big(\ \mathrm{Obj}=\{\mathcal A\},\;
\mathrm{Mor}=\{\alpha_{\mathrm{ext}}(\Lambda)\mid \Lambda\in G_{\mathrm{ext}}\}\ \big),
\]
die das geometrische Hintergrundsystem beschreibt.
Alle kovarianten Objekte (Felder, Operatoren, Gleichungen) sind Funktoren
\[
\mathbf{Back}(\mathcal A)\longrightarrow\mathbf{C}
\]
in eine geeignete Zielkategorie $\mathbf{C}$ (z.\,B.\ $\mathbf{Bund}$, $\mathbf{Set}$, $\mathbf{Vec}$).

\bigskip
\emph{2. Innere Symmetrien und Bündelstruktur.}

\textbf{(a) Strukturgruppe.}
Es sei $G_{\mathrm{int}}$ eine Lie-Gruppe, die als \emph{innere Symmetriegruppe}
fungiert. Ihr geometrisches Trägerobjekt ist ein Hauptbündel
\[
\pi_P:P\longrightarrow \mathcal A
\]
mit freier, glatter Rechtsaktion $R_g:P\to P$, $(p,g)\mapsto p\cdot g$.

\textbf{(b) Kategorie der Hauptbündel.}
Die Objekte sind Paare $(P\to\mathcal A,G_{\mathrm{int}})$,
die Morphismen $f:P_1\to P_2$ sind glatte, $G_{\mathrm{int}}$–äquivariante
Bündelabbildungen über der Identität $\mathrm{id}_\mathcal A$.
\[
\mathbf{Prin}_{G_{\mathrm{int}}}(\mathcal A)
:=\big\{P\to\mathcal A\text{ Hauptbündel mit Strukturgruppe }G_{\mathrm{int}}\big\}.
\]

\textbf{(c) Assoziierte Bündel.}
Für jede Darstellung $\rho_{\mathrm{int}}:G_{\mathrm{int}}\to \mathrm{GL}(F_0)$
definiere das \emph{assoziierte Bündel}
\[
\pi:\ \mathcal F := P\times_{\rho_{\mathrm{int}}}F_0 \longrightarrow \mathcal A,
\]
als den Quotienten $(P\times F_0)/{\sim}$ mit $(p\cdot g, f)\sim(p,\rho_{\mathrm{int}}(g)f)$.

\textbf{(d) Lokale Trivialisierungen.}
Für offene Mengen $U\subset \mathcal A$ und lokale Schnitte $s:U\to P$ ergibt sich
eine lokale Darstellung
\[
\mathcal F|_U \cong U\times F_0,\quad
[(p,f)]\mapsto (\pi(p), \rho_{\mathrm{int}}(g(s(\pi(p)),p))f).
\]
Diese Darstellung definiert eine Garbe $\mathfrak F$ lokaler Abschnitte
$U\mapsto\Gamma(U,\mathcal F|_U)$.

\bigskip
\emph{3. Feldfunktor und Garbenstruktur.}

\textbf{(a) Definition.}
Der \emph{Feldfunktor} eines Modells ist der Garben-Funktor
\[
\mathfrak F:\ \mathrm{Open}(\mathcal A)^{\mathrm{op}}\longrightarrow\mathbf{Set},
\qquad
U\longmapsto \mathfrak F(U):=\Gamma(U,\mathcal F|_U),
\]
mit Restriktionsabbildungen $r_{U,V}$ für $U\subseteq V$.

\textbf{(b) Garbenaxiome.}
Für eine Überdeckung $U=\bigcup_i U_i$ gilt:
\begin{enumerate}
\item (\emph{Konsistenz}) Wenn $\phi_i\in\mathfrak F(U_i)$ mit
$r_{U_i\cap U_j,U_i}\phi_i=r_{U_i\cap U_j,U_j}\phi_j$ für alle $i,j$,
dann existiert $\phi\in\mathfrak F(U)$ mit $r_{U_i,U}\phi=\phi_i$.
\item (\emph{Eindeutigkeit}) Wenn $\phi,\psi\in\mathfrak F(U)$ mit
$r_{U_i,U}\phi=r_{U_i,U}\psi$ für alle $i$, dann $\phi=\psi$.
\end{enumerate}
Damit ist $\mathfrak F$ eine (prä)Garbenstruktur im Sinn der lokalen Feldbeschreibung.

\textbf{(c) Lokale und globale Beschreibungen.}
Ein \emph{globale Beschreibung} ist $\phi\in\mathfrak F(\mathcal A)$,
eine \emph{lokale Beschreibung} ein Paar $(U,\phi)$ mit $\phi\in\mathfrak F(U)$.
Die Gesamtheit aller lokalen Beschreibungen bildet die kolimitartige Familie
$\mathrm{colim}_{U\subset\mathcal A}\mathfrak F(U)$.

\bigskip
\emph{4. Wirkungen der Symmetrien auf Feldern.}

\textbf{(a) Externe Wirkung.}
Die Aktion $G_{\mathrm{ext}}\curvearrowright\mathcal A$ induziert auf $\mathcal F$
eine \emph{kovariante} Wirkung durch Bündelautomorphismen
\[
\rho_{\mathrm{ext}}(\Lambda):\mathcal F\to \mathcal F,
\quad \pi\circ\rho_{\mathrm{ext}}(\Lambda)=\alpha_{\mathrm{ext}}(\Lambda)\circ\pi,
\]
die die Faserstruktur respektieren.  
Die induzierte Wirkung auf Abschnitten:
\[
(\rho_{\mathrm{ext}}(\Lambda)\cdot \phi)(x)
:=\rho_{\mathrm{ext}}(\Lambda)\big(\phi(\Lambda^{-1}\cdot x)\big),
\qquad \phi\in\Gamma(\mathcal F).
\]

\textbf{(b) Innere Wirkung.}
Die Eichgruppe $\mathcal G_{\mathrm{int}}(U)$ wirkt \emph{vertikal} auf
$\mathcal F|_U$:
\[
(\rho_{\mathrm{int}}(\gamma)\cdot \phi)(x)
:=\rho_{\mathrm{int}}\big(\gamma(x)\big)\,\phi(x),
\qquad \gamma\in\mathcal G_{\mathrm{int}}(U).
\]
Für eine globale Darstellung $\rho_{\mathrm{int}}:G_{\mathrm{int}}\to GL(F_0)$
reduziert sich dies zu einer Punkt-für-Punkt-Wirkung auf jeder Faser.

\textbf{(c) Kombinierte Wirkung.}
Die kombinierte Gruppe $G_{\mathrm{tot}}=G_{\mathrm{ext}}\ltimes\mathcal G_{\mathrm{int}}(-)$
wirkt auf den Feldfunktor $\mathfrak F$ als Garben-Automorphismengruppe:
\[
(\Lambda,\gamma)\cdot\phi
:=\rho_{\mathrm{ext}}(\Lambda)\big(\rho_{\mathrm{int}}(\gamma)\phi\big).
\]
Formal: $\mathfrak F$ ist ein $G_{\mathrm{tot}}$–äquivarianter Funktor
\[
\mathfrak F: (\mathrm{Open}(\mathcal A)^{\mathrm{op}}\times G_{\mathrm{tot}})\longrightarrow\mathbf{Set}.
\]

\bigskip
\emph{5. Gruppoid der Beschreibungen.}

\textbf{(a) Objekte und Morphismen.}
Definiere das \emph{Beschreibungsgrouppoid}
\[
\mathsf{Desc}(\mathsf{Sys})
:=\Big(
\mathrm{Obj}=\{(U,\phi)\mid U\subset\mathcal A,\ \phi\in\mathfrak F(U)\},\quad
\mathrm{Mor}=\{(\Lambda,\gamma):(U,\phi)\to(U',\phi')\}\Big),
\]
wobei $(\Lambda,\gamma)$ genau dann ein Morphismus ist, wenn
\[
\alpha_{\mathrm{ext}}(\Lambda)(U)=U',\qquad
\rho_{\mathrm{ext}}(\Lambda)\rho_{\mathrm{int}}(\gamma)\phi=\phi'\ \text{auf}\ U\cap U'.
\]
Komposition ist punktweise definiert:
\[
(\Lambda_2,\gamma_2)\circ(\Lambda_1,\gamma_1)
=(\Lambda_2\Lambda_1,\,\gamma_2\,\rho_{\mathrm{int}}(\Lambda_2)\gamma_1),
\]
Identität $(\mathrm{id}_{\mathcal A},e)$.

\textbf{(b) Struktur.}
$\mathsf{Desc}(\mathsf{Sys})$ ist eine interne Kategorie in $\mathbf{Set}$
(bzw.\ ein Gruppoid in Garben über $\mathcal A$).
\begin{itemize}
\item Jede Beschreibung $(U,\phi)$ ist ein Objekt.
\item Jeder zulässige Transformationswechsel ist ein Isomorphismus.
\item Die Stabilisatorgruppe eines $\phi$ ist
\[
\mathrm{Stab}(\phi)
=\{(\Lambda,\gamma)\in G_{\mathrm{tot}}(U)\mid
 (\Lambda,\gamma)\cdot\phi=\phi\}.
\]
\end{itemize}

\textbf{(c) Orbitraum und Zustände.}
Der \emph{Orbit} eines $\phi\in\mathfrak F(U)$ ist
$G_{\mathrm{tot}}(U)\cdot\phi$.  
Die Menge der Orbits bildet den \emph{Zustandsraum}
\[
\mathcal{S}tate(\mathsf{Sys})
:= \Gamma(\mathcal F)\big/\big(G_{\mathrm{ext}}\times\mathcal G_{\mathrm{int}}(\mathcal A)\big),
\]
d.\,h.\ die Raumzeit- und Eichäquivalenzklassen der Feldkonfigurationen.

\textbf{(d) Funktorielle Sicht.}
$\mathsf{Desc}(\mathsf{Sys})$ ist ein Gruppoid-Objekt in der 2-Kategorie der
Garben über $\mathcal A$,
\[
\mathsf{Desc}(\mathsf{Sys}) \in \mathbf{Gpd}(\mathbf{Sh}(\mathcal A)).
\]
In dieser Sprache wird der gesamte Raum der Beschreibungen mitsamt den
Wechseln (Bezugssystemen, Eichtransformationen) als einheitliches geometrisches
Objekt modelliert. Physikalisch relevante Objekte sind Funktoren oder
natürliche Transformationen auf diesem Gruppoid.


\emph{6. Natürliche Differentialoperatoren und Kovarianz.}

\textbf{(a) Differentialoperatoren als natürliche Transformationen.}
Sei $J^k\mathcal F$ das $k$-Jet-Bündel des Bündels $\mathcal F\to\mathcal A$.
Ein \emph{lokaler Differentialoperator der Ordnung $\le k$} von $\mathcal F$
nach einem Zielbündel $\mathcal E\to\mathcal A$ ist eine Bündelabbildung
\[
D:\ J^k\mathcal F \longrightarrow \mathcal E,
\]
die glatte Abbildungen zwischen den jeweiligen Totalräumen induziert und sich zu einem
Funktor auf offenen Teilmengen fortsetzt:
\[
D_U:\ J^k(\mathfrak F(U)) \longrightarrow \mathfrak E(U),
\qquad
\mathfrak E(U):=\Gamma(U,\mathcal E|_U).
\]
$D$ ist \emph{natürlich} relativ zur Kategorie $\mathbf{Back}(\mathcal A)$,
wenn für alle Morphismen $\Lambda:\mathcal A\to\mathcal A$ gilt:
\[
\rho_{\mathrm{ext}}(\Lambda)\circ D_U
= D_{\Lambda U}\circ J^k(\rho_{\mathrm{ext}}(\Lambda)),
\]
und analog für innere Transformationen $\gamma$.

\textbf{(b) Lokale Ausdrücke.}
In lokalen Koordinaten $(x^i,u^\alpha,u^\alpha_I)$ des Jetbündels
($I$ multiindex) besitzt $D$ die Form
\[
(D\phi)^\beta(x)
= F^\beta\big(x,u^\alpha(x),\partial_I u^\alpha(x)\big)
\]
für glatte Funktionen $F^\beta$.
Natürliche Operatoren sind genau diejenigen, deren $F^\beta$ tensorielle
Verhaltensregeln erfüllen, die aus der Geometrie von $\mathcal A$ folgen.

\textbf{(c) Gleichungsmenge als Nullstellenfunktor.}
Definiere für jedes offene $U\subseteq\mathcal A$
\[
\mathcal E(U) := \{\phi\in\mathfrak F(U)\mid D_U(j^k\phi)=0\}.
\]
Dies definiert einen Unterfunktor $\mathcal E\subset\mathfrak F$,
den \emph{Gleichungsfunktor}.
Kovarianz von $\mathcal E$ folgt unmittelbar aus Natürlichkeit von $D$:
\[
D_{\Lambda U}(j^k((\Lambda,\gamma)\cdot\phi))
= \rho_{\mathrm{ext}}(\Lambda)\rho_{\mathrm{int}}(\gamma)
\big(D_U(j^k\phi)\big),
\]
woraus $(\Lambda,\gamma)\cdot\mathcal E(U)=\mathcal E(\Lambda U)$ folgt.

\textbf{(d) Äquivarianzstruktur.}
$D$ ist damit ein $G_{\mathrm{tot}}$-äquivarianter Morphismus in der Kategorie
der Bündel über $\mathcal A$,
\[
D:\ J^k\mathcal F \longrightarrow \mathcal E,\qquad
D\in \mathrm{Hom}_{\mathbf{Bund}_{G_{\mathrm{tot}}}}(J^k\mathcal F,\mathcal E).
\]

\textbf{(e) Kommutative Diagrammform.}
\[
\begin{array}{ccc}
J^k\mathfrak F(U) & \xrightarrow{D_U} & \mathfrak E(U)\\
\downarrow J^k(\rho(\Lambda,\gamma)) & & \downarrow \rho(\Lambda,\gamma)\\
J^k\mathfrak F(\Lambda U) & \xrightarrow{D_{\Lambda U}} & \mathfrak E(\Lambda U)
\end{array}
\quad
\Rightarrow\quad
\mathcal E(U)=\ker D_U\ \text{ ist kovariant.}
\]

\bigskip
\emph{7. Variations- und Lagrangeformalismus.}

\textbf{(a) Lagrangedichten.}
Eine \emph{Lagrangedichte der Ordnung $\le k$} ist eine horizontale $n$-Form
\[
\mathcal L\in \Omega^{n,0}(J^k\mathcal F),
\qquad n=\dim\mathcal A,
\]
die als natürliches Objekt unter $\mathbf{Back}(\mathcal A)$ transformiert,
d.\,h.
\[
(\Lambda,\gamma)^\ast\mathcal L = \mathcal L
\]
bis auf eine total-derivative Form $d_h\Theta$.
Dies garantiert Kovarianz der entstehenden Euler–Lagrange-Gleichungen.

\textbf{(b) Variationsbikomplex.}
Im Variationsbikomplex auf $J^k\mathcal F$ werden die Differentiale
$d=d_h+d_v$ in horizontale und vertikale Komponenten zerlegt.
Die Euler–Lagrange-Form $EL(\mathcal L)$ ist durch
\[
d_v\mathcal L = EL(\mathcal L) + d_h\Theta
\]
definiert; sie bestimmt den natürlichen Differentialoperator
$EL(\mathcal L):J^{2k}\mathcal F\to V^\ast(\mathcal F)\otimes\Lambda^nT^\ast\mathcal A$,
wobei $V^\ast(\mathcal F)$ das duale Vertikalbündel ist.

\textbf{(c) Gleichungen aus der Variation.}
\[
\mathcal E_{\mathcal L}(U)
=\{\phi\in\mathfrak F(U)\mid EL(\mathcal L)_U(j^{2k}\phi)=0\}.
\]
$EL(\mathcal L)$ ist ein natürlicher Differentialoperator,
daher ist $\mathcal E_{\mathcal L}$ automatisch $G_{\mathrm{tot}}$-kovariant.

\textbf{(d) Noether-1-Struktur.}
Jede einparametrige Familie $(\Lambda_t,\gamma_t)\in G_{\mathrm{tot}}$
mit infinitesimalem Generator $\xi$ und
$\mathcal L$-Symmetrieeigenschaft
$\mathcal L((\Lambda_t,\gamma_t)\cdot\phi)=\mathcal L(\phi)+d_h\Theta_\xi$
liefert einen (lokal definierten) Strom $J_\xi$ mit
$d_h J_\xi = EL(\mathcal L)\cdot \delta_\xi\phi$.
Auf der Lösungsmenge $EL(\mathcal L)=0$ gilt $d_h J_\xi=0$,
d.\,h. lokale Erhaltungsgesetze sind exakt
die homologische Folge natürlicher Symmetrien der Lagrange-Garbenstruktur.

\bigskip
\emph{8. Funktor der Lösungen und Observablen.}

\textbf{(a) Lösungsfunktor.}
Der Lösungsfunktor ist
\[
\mathcal S ol:\ \mathrm{Open}(\mathcal A)^{\mathrm{op}}\longrightarrow \mathbf{Set},
\quad
U\longmapsto \mathcal E(U).
\]
Er ist ein Unterfunktor von $\mathfrak F$ und $G_{\mathrm{tot}}$-äquivariant:
\[
(\Lambda,\gamma)\cdot \mathcal S ol(U)
= \mathcal S ol(\Lambda U).
\]

\textbf{(b) Natur und Monoidalstruktur.}
$\mathcal S ol$ kann als Objekt einer monoidalen Kategorie betrachtet werden,
wenn $\mathfrak F$ eine additive oder tensorielle Struktur trägt
(z.\,B. direkte Summe oder Produkt von Feldern).
Kovarianz entspricht dann der Monoidalkompatibilität
\[
\mathcal S ol(\Lambda U)\cong\rho(\Lambda)\mathcal S ol(U).
\]

\textbf{(c) Observablen als natürliche Transformationen.}
Eine \emph{Observable} ist eine natürliche Transformation
\[
O:\ \mathcal S ol \Rightarrow \underline{\mathbb K},
\]
wobei $\underline{\mathbb K}$ der konstante Funktor auf einem Körper
$\mathbb K$ ist.  
$O$ ist $G_{\mathrm{tot}}$-invariant, wenn
\[
O_U = O_{\Lambda U}\circ (\Lambda,\gamma),
\quad \forall (\Lambda,\gamma).
\]
Die Familie aller Observablen bildet die Kategorie
\[
\mathcal O := \mathrm{Nat}(\mathcal S ol,\underline{\mathbb K}).
\]

\textbf{(d) Algebraische Strukturen.}
Falls $\mathcal S ol(U)$ eine algebraische Struktur trägt
(z.\,B. Poisson-Struktur, symplektische Form),
induzieren die Observablen durch Auswertung
eine kommutative oder *-Algebra
\[
\mathcal O(U)=\mathrm{Fun}_{\mathrm{alg}}(\mathcal S ol(U),\mathbb K),
\]
die wieder zu einem Garben- oder Prägarben-Objekt wird.

\bigskip
\emph{9. Quotienten und Stack-Struktur.}

\textbf{(a) Quotienten-Gruppoid.}
Da die $G_{\mathrm{tot}}$-Wirkung typischerweise nicht frei ist,
ist der Raum physikalischer Zustände kein Mengenquotient, sondern
ein Gruppoid- oder Stack-Quotient:
\[
[\mathcal S ol / G_{\mathrm{tot}}]
\in \mathbf{Gpd}(\mathbf{Sh}(\mathcal A)).
\]
Objekte sind Lösungen, Morphismen Eich-/Bezugstransformationen.

\textbf{(b) Invariante Observablen.}
Physikalische Observablen sind Funktoren auf diesem Stack:
\[
\mathcal O_{\mathrm{phys}}
=\mathrm{Funct}\big([\mathcal S ol / G_{\mathrm{tot}}],\,\mathbf{Set}\big)
\simeq \mathrm{Nat}\big(\mathcal S ol,\underline{\mathbb K}\big)^{G_{\mathrm{tot}}}.
\]
Invarianz bedeutet Funktorialität über den Quotienten.

\bigskip
\emph{10. Zusammenfassende Diagrammatische Struktur.}

\[
\begin{array}{ccccc}
G_{\mathrm{tot}}
& \curvearrowright & \mathfrak F & \xrightarrow{D} & \mathfrak E\\[0.4em]
& & \rotatebox{90}{$\subset$} & & \\[-0.1em]
& & \mathcal S ol=\ker D & &
\end{array}
\qquad
\begin{array}{l}
\text{Kovarianz: } D\text{ natürlich}\\[0.3em]
\text{Invarianz: } O\in\mathrm{Nat}(\mathcal S ol,\underline{\mathbb K})
\end{array}
\]

Alle Gleichungen, Operatoren und Observablen sind damit als
natürliche oder äquivariante Transformationen zwischen Funktoren modelliert.
Invarianz physikalischer Gleichungen entspricht der
Kommutativität solcher Diagramme; die Kategorie der möglichen
Beschreibungen (Bezugssysteme und Eichungen) bildet die
Struktur, über der diese Natürlichkeit definiert ist.

\pagebreak


\subsection{Beispiel: Freies Dirac-Feld auf dem Minkowski-Raum}

\emph{Ziel.}
Wir instanziieren den abstrakten Rahmen
\[
\mathsf{Sys}=(\mathcal A,\,G_{\mathrm{ext}},\,G_{\mathrm{int}},\,\mathcal F,\,\rho_{\mathrm{ext}},\,\rho_{\mathrm{int}},\,\mathcal E,\,\mathcal O)
\]
für das freie Dirac-Feld in spezieller Relativität (ohne Eichkopplung).

\medskip
\emph{1. Arena \(\mathcal A\).}
\[
\mathcal A=(M,g),\qquad M=\R^{1,3},\qquad g=\mathrm{diag}(1,-1,-1,-1).
\]
Die Kategorie \(\mathbf{Geom}\) sei „pseudoriemannsche Mannigfaltigkeiten mit Isometrien“.

\medskip
\emph{2. Äußere Symmetrie \(G_{\mathrm{ext}}\).}
\[
G_{\mathrm{ext}}=\mathcal P:=\R^{1,3}\rtimes SO^+(1,3)
\]
(Poincaré-Gruppe). Wirkung auf der Arena:
\[
\alpha_{\mathrm{ext}}(a,\Lambda):\ x\mapsto \Lambda x + a.
\]
Isometrien: \(\Lambda^T g\,\Lambda=g\), \(\det\Lambda=1\), \(\Lambda^0{}_0>0\).

\medskip
\emph{3. Innere Symmetrie \(G_{\mathrm{int}}\).}
Im freien Dirac-Modell ohne Ladung setzen wir
\[
G_{\mathrm{int}}=\{e\}\quad(\text{trivial}).
\]
(Bei QED wäre \(G_{\mathrm{int}}=U(1)\), siehe Hinweis unten.)

\medskip
\emph{4. Feldbündel \(\mathcal F\).}
Spin-Struktur auf \((M,g)\) ist trivial; das zugehörige Spinor\,-\,Bündel ist global trivial:
\[
\pi:\ \mathcal F=M\times \Sigma \longrightarrow M,\qquad \Sigma=\C^4.
\]
Abschnitte \(\psi\in\Gamma(\mathcal F)\cong C^\infty(M,\C^4)\) heißen Dirac-Spinorfelder.

\medskip
\emph{5. Darstellungen \(\rho_{\mathrm{ext}},\rho_{\mathrm{int}}\).}
\[
\rho_{\mathrm{ext}}|_{SO^+(1,3)}=\sigma:\ \Spin(1,3)\to GL(\Sigma),
\quad \sigma\ \text{kompatibel mit}\ \gamma^\mu:
\]
\[
\sigma(a)^{-1}\,\gamma^\mu\,\sigma(a)=\Lambda^\mu{}_\nu\,\gamma^\nu,\qquad \Lambda=\rho(a)\in SO^+(1,3).
\]
Auf Translationen wirkt \(\rho_{\mathrm{ext}}(a,\mathbf 1)\) durch Pullback:
\((\rho_{\mathrm{ext}}(a,\mathbf 1)\psi)(x)=\psi(x-a)\).
Da \(G_{\mathrm{int}}\) trivial ist, gilt \(\rho_{\mathrm{int}}=\mathrm{id}\).

\medskip
\emph{6. Feldfunktor und Beschreibungen.}
Für offene \(U\subset M\):
\[
\mathfrak F(U):=\Gamma(U,\mathcal F|_U)\cong C^\infty(U,\C^4),\qquad
r_{U,V}:\phi\mapsto \phi|_U.
\]
Eine \emph{Beschreibung} ist \((U,\phi)\) mit \(\phi\in \mathfrak F(U)\).

\medskip
\emph{7. Natürlicher Differentialoperator \(D\).}
Wähle Gamma-Matrizen \(\gamma^\mu\in \End(\Sigma)\) mit
\(\{\gamma^\mu,\gamma^\nu\}=2g^{\mu\nu}\mathbf 1\).
Definiere den natürlichen Operator erster Ordnung
\[
D:\ J^1\mathcal F\to \mathcal E\ (=\mathcal F),\qquad
(D\psi)(x):= i\,\gamma^\mu\,\partial_\mu \psi(x)-m\,\psi(x).
\]
Natürlichkeit (Kovarianz) heißt:
\[
\rho_{\mathrm{ext}}(a,\Lambda)\circ D
= D\circ J^1(\rho_{\mathrm{ext}}(a,\Lambda)),
\]
was wegen \(\sigma(a)^{-1}\gamma^\mu\sigma(a)=\Lambda^\mu{}_\nu\gamma^\nu\) erfüllt ist.

\medskip
\emph{8. Gleichungsmenge \(\mathcal E\).}
Für \(U\subset M\):
\[
\mathcal E(U):=\{\psi\in\mathfrak F(U)\mid (i\gamma^\mu\partial_\mu-m)\psi=0\ \text{auf }U\}.
\]
Das ist die (lokale) Lösungsmenge der Dirac-Gleichung. Kovarianz:
\((a,\Lambda)\cdot \mathcal E(U)=\mathcal E(\Lambda U+a)\).

\medskip
\emph{9. Gruppoid der Beschreibungen.}
Objekte \((U,\psi)\), Morphismus \((a,\Lambda):(U,\psi)\to(U',\psi')\) genau dann, wenn
\[
U'=\Lambda U + a,\qquad
\psi'(x')=S(\Lambda)\,\psi(\Lambda^{-1}(x'-a)),\quad S(\Lambda)=\sigma(a)\in \Spin(1,3).
\]
Alle Morphismen sind Isomorphismen; Orbits sind Bezugssystemäquivalenzklassen von Lösungen.

\medskip
\emph{10. Observablen \(\mathcal O\).}
Beispiel eines natürlichen (lokalen) Funktionals:
\[
j^\mu(\psi):=\bar\psi\,\gamma^\mu\,\psi,\qquad \bar\psi:=\psi^\dagger\gamma^0,
\]
mit \(\partial_\mu j^\mu=0\) auf \(\mathcal E\).
Globale Observablen entstehen durch Integration (sofern wohldefiniert):
\[
Q(\psi):=\int_\Sigma j^\mu n_\mu\, d\Sigma
\]
(unabhängig von \(\Sigma\) für geeignete Abklingbedingungen). Invarianz: \(Q\) ist \(G_{\mathrm{ext}}\)-invariant.

\medskip
\emph{11. Diagrammatische Kovarianz.}
\[
\begin{array}{ccc}
J^1\mathfrak F(U) & \xrightarrow{\ \ D_U\ \ } & \mathfrak F(U)\\[0.3em]
\downarrow J^1(\rho_{\mathrm{ext}}(a,\Lambda)) & & \downarrow \rho_{\mathrm{ext}}(a,\Lambda)\\[0.3em]
J^1\mathfrak F(\Lambda U{+}a) & \xrightarrow{\ \ D_{\Lambda U{+}a}\ \ } & \mathfrak F(\Lambda U{+}a)
\end{array}
\qquad\Rightarrow\qquad
\mathcal E(U)=\ker D_U\ \text{ ist kovariant.}
\]

\medskip
\emph{12. Spezifikation der \(\gamma\)-Struktur.}
Wähle etwa Dirac-Darstellung:
\[
\gamma^0=\begin{pmatrix} I & 0\\ 0 & -I\end{pmatrix},\qquad
\gamma^i=\begin{pmatrix} 0 & \sigma^i\\ -\sigma^i & 0\end{pmatrix},\quad
\gamma^5:=i\gamma^0\gamma^1\gamma^2\gamma^3.
\]
Dann \(\Sigma=\C^4\), \(\sigma^{\mu\nu}:=\tfrac{i}{2}[\gamma^\mu,\gamma^\nu]\),
\(S(\Lambda)=\exp\!\big(-\tfrac{i}{4}\omega_{\mu\nu}\sigma^{\mu\nu}\big)\).

\medskip
\emph{13. Zusammenfassung der Identifikationen.}
\[
\begin{array}{ll}
\text{Arena} & \mathcal A=(\R^{1,3},g)\\
\text{äußere Gruppe} & G_{\mathrm{ext}}=\R^{1,3}\rtimes SO^+(1,3)\\
\text{innere Gruppe} & G_{\mathrm{int}}=\{e\}\\
\text{Feldbündel} & \mathcal F=M\times \C^4\\
\text{Darstellungen} & \rho_{\mathrm{ext}}:\ \text{Poincaré}\to \Aut(\mathcal F),\ \rho_{\mathrm{int}}=\mathrm{id}\\
\text{Operator} & D=i\gamma^\mu\partial_\mu - m\\
\text{Gleichungen} & \mathcal E=\{\psi\mid D\psi=0\}\\
\text{Observablen} & \mathcal O:\ \psi\mapsto \text{invariante Funktionale (z.\,B.\ Ladung)}\\
\text{Gruppoid} & \mathsf{Desc}:\ \text{alle }(U,\psi)\ \text{und Poincaré-Wechsel}
\end{array}
\]

\medskip
\emph{Hinweis (Eichkopplung/QED).}
Für QED ersetzt man \(G_{\mathrm{int}}=\{e\}\) durch \(U(1)\), wählt ein Hauptbündel
\(P\to M\), das assoziierte Spinorbündel \(\mathcal F\) und den Anschluss \(A\) auf \(P\).
Der natürliche Operator wird \(D=i\gamma^\mu(\partial_\mu + iq A_\mu) - m\),
die innere Wirkung \(\rho_{\mathrm{int}}(e^{i\alpha}):\ \psi\mapsto e^{-iq\alpha}\psi\),
\(A\mapsto A+d\alpha\). Alle obigen Punkte bleiben formal gleich, mit zusätzlicher
\(U(1)\)-Äquivarianz.


\pagebreak

\subsection{Konzeptuelle Analyse: Warum äußere und innere Symmetrien als Gruppenaktionen auftreten}

\emph{Ausgangspunkt.}
Ein physikalisches System wird durch eine Menge von Beschreibungen (Felder, Koordinaten, Bezugssysteme)
repräsentiert.  
Zwischen diesen Beschreibungen existieren „Wechsel“, die die gleiche physikalische Situation
in verschiedenen Darstellungsformen ausdrücken.
Die Struktur dieser Wechsel erzwingt, dass sie algebraisch durch Gruppen- oder Gruppoidaktionen
modelliert werden.  

Am Beispiel des freien Dirac-Feldes im Minkowski-Raum lässt sich dies formal nachvollziehen.

\bigskip
\emph{1. Kinematische Arena als geometrische Kategorie.}

Die Raumzeit selbst ist ein Objekt einer geometrischen Kategorie:
\[
\mathcal A=(M,g)\in\mathbf{Geom},\qquad M=\R^{1,3},\ g=\mathrm{diag}(1,-1,-1,-1).
\]
Jede Transformation, die den geometrischen Aufbau (Metrikstruktur) erhält, ist ein
Automorphismus in dieser Kategorie.
Die Gesamtheit dieser Automorphismen bildet eine Lie-Gruppe:
\[
\Aut_{\mathbf{Geom}}(\mathcal A)=\mathrm{Isom}(M,g)=\mathcal P
=\R^{1,3}\rtimes SO^+(1,3).
\]
\textbf{Erkenntnis:}
Jede physikalische Theorie auf $(M,g)$, deren Gleichungen die Metrikstruktur respektieren,
muss invariant unter diesen Isometrien sein.
Dadurch wird \emph{äußere Symmetrie} notwendigerweise durch Gruppenaktionen beschrieben.

\bigskip
\emph{2. Äußere Symmetrien als Automorphismen der Arena.}

Eine äußere Transformation ist eine Abbildung der Arena auf sich selbst:
\[
\alpha_{\mathrm{ext}}(\Lambda,a):\ x\mapsto \Lambda x + a.
\]
Sie ist invertierbar und die Komposition ist wieder vom gleichen Typ:
\[
\alpha_{\mathrm{ext}}(\Lambda_2,a_2)\circ\alpha_{\mathrm{ext}}(\Lambda_1,a_1)
=\alpha_{\mathrm{ext}}(\Lambda_2\Lambda_1,\Lambda_2 a_1+a_2).
\]
\textbf{Erkenntnis:}
Die Menge aller solcher Abbildungen bildet algebraisch eine Gruppe
— die Poincaré-Gruppe —, und diese Gruppenstruktur entspricht genau der Kompositionsstruktur
von Koordinatenwechseln.  
Die äußere Symmetriegruppe ist also nichts anderes als die interne Automorphismengruppe
der geometrischen Arena.

\bigskip
\emph{3. Warum Felder als Bündelobjekte modelliert werden.}

Ein Feld ist lokal eine Funktion \(x\mapsto \psi(x)\), deren Werte nicht in \(\R\),
sondern in einem linearen Raum \(F_0\) liegen (z.\,B. $\C^4$ für Spinoren).  
Um lokale Koordinatenwechsel konsistent zu formulieren, muss das Feld unter
einer Transformation der Basis (Bezugssystemwechsel) eine lineare Abbildung
auf seinen Werten induzieren:
\[
x\mapsto \Lambda x\ \Rightarrow\
\psi'(x')=S(\Lambda)\psi(x),
\]
wobei \(S(\Lambda)\) eine lineare Automorphismusgruppe des Faserraums $F_0$ ist.
Damit wird jedes Feld formal ein Schnitt eines Faserbündels
\[
\pi:\ \mathcal F = M\times F_0 \to M.
\]
\textbf{Erkenntnis:}
Damit äußere Transformationen die gleiche physikalische Feldstruktur erhalten,
muss ihre Wirkung durch lineare Automorphismen auf den Fasern vermittelt werden.
Dies ist die Definition einer \emph{Darstellung der Symmetriegruppe} auf dem Faserraum.  

Das Bündel wird so zu einem \emph{repräsentierenden Objekt} der externen Symmetrie.

\bigskip
\emph{4. Warum interne Symmetrien als vertikale Automorphismen erscheinen.}

Physikalische Theorien enthalten oft Freiheitsgrade, die nicht durch Raumzeitbewegung,
sondern durch \emph{interne Wahl von Bezugssystemen in der Faser} beschrieben werden:
die Eichfreiheit.

Eine Eichtransformation ändert nur die lokale Faserbasis,
nicht den Basispunkt:
\[
(\rho_{\mathrm{int}}(\gamma)\psi)(x)
=\rho_{\mathrm{int}}(\gamma(x))\,\psi(x),
\quad \gamma(x)\in G_{\mathrm{int}}.
\]
Dies sind genau die \emph{vertikalen Automorphismen} des Faserbündels $\mathcal F$,
da sie die Projektion $\pi$ unverändert lassen:
\[
\pi\circ\rho_{\mathrm{int}}(\gamma)=\pi.
\]
\textbf{Erkenntnis:}
Die Forderung, dass man lokal eine Faserbasis frei wählen darf (lokale Eichfreiheit),
impliziert strukturell, dass die Gesamtheit solcher Vertikaltransformationen eine
Lie-Gruppe bilden muss, die faserweise wirkt — also eine
Garbenstruktur $U\mapsto\mathcal G_{\mathrm{int}}(U)$.

\bigskip
\emph{5. Vereinheitlichung durch Funktorialität.}

Sowohl äußere als auch innere Symmetrien wirken funktoriell:
\begin{align*}
\rho_{\mathrm{ext}}:&\ G_{\mathrm{ext}}\to \Aut_{\mathbf{Bund}}(\mathcal F),
&
(\Lambda,a)\mapsto\rho_{\mathrm{ext}}(\Lambda,a),\\[0.3em]
\rho_{\mathrm{int}}:&\ G_{\mathrm{int}}\to \Aut_{\mathbf{Bund}}(\mathcal F),
&
\gamma\mapsto\rho_{\mathrm{int}}(\gamma).
\end{align*}
Da beide Abbildungen Gruppenhomomorphismen sind, folgen sie denselben
Kompositionsregeln wie ihre Ursprungsgruppen.
Dies gewährleistet, dass hintereinander ausgeführte Transformationen
wieder eine Transformation gleicher Art ergeben — genau die Eigenschaft,
die physikalisch \emph{Konsistenz von Koordinatenwechseln und Eichwechseln}
ausdrückt.

\bigskip
\emph{6. Äquivariante Garbenstruktur als formaler Ausdruck der Invarianz.}

Die Aktion $(\Lambda,\gamma)$ auf der Garbe der Felder
\[
\mathfrak F(U)=\Gamma(U,\mathcal F|_U)
\]
wird als natürliche Transformation von Funktoren formuliert:
\[
\rho(\Lambda,\gamma):\ \mathfrak F(U)\longrightarrow\mathfrak F(\Lambda U),
\quad
\psi\mapsto\rho_{\mathrm{ext}}(\Lambda)\rho_{\mathrm{int}}(\gamma)\psi.
\]
Die Natürlichkeit (Kommutativität der Diagramme)
\[
\rho(\Lambda_2,\gamma_2)\circ\rho(\Lambda_1,\gamma_1)
=\rho((\Lambda_2,\gamma_2)\circ(\Lambda_1,\gamma_1))
\]
ist formal die Kategorieversion der Forderung:
„Wechsel von Beschreibungen ist assoziativ und invertierbar.“
Das macht $\mathfrak F$ zu einem $G_{\mathrm{tot}}$–äquivarianten Funktor.

\bigskip
\emph{7. Warum dies physikalisch unvermeidlich ist.}

\textbf{(i) Beobachterwechsel (äußere Symmetrie).}
Dasselbe physikalische Ereignis wird in verschiedenen Koordinatensystemen
unterschiedlich beschrieben.  
Damit „dieselbe Physik“ vorliegt, muss es eine konsistente Übersetzung
zwischen Beschreibungen geben — formal ein Isomorphismus.
Daher ist der Raum der möglichen Beschreibungen ein \emph{Gruppoid}:
Objekte sind lokale Darstellungen, Morphismen die Wechsel.

\textbf{(ii) Eichfreiheit (innere Symmetrie).}
Auch ohne Wechsel des Beobachters kann man lokal die Faserbasis ändern,
ohne die Physik zu ändern.
Diese zusätzliche Freiheit ist wieder ein System von Isomorphismen,
das über jeder offenen Menge wirkt — also ein
„faserweises Gruppoid in Garben“ (die Eichgarbe).

\textbf{(iii) Strukturelle Notwendigkeit.}
Sobald man fordert:
\begin{itemize}
\item die Existenz mehrerer äquivalenter Beschreibungen eines Zustands,
\item deren Zusammensetzung wieder eine Beschreibung liefert,
\item und die Umkehrbarkeit jedes Wechsels,
\end{itemize}
ergibt sich zwingend die Gruppoidstruktur.
Die klassischen „Symmetriegruppen“ sind ihre Automorphismengruppen.

\bigskip
\emph{8. Anwendung auf das Dirac-Feld.}

Im Dirac-Beispiel konkretisiert sich dies zu:
\[
\begin{array}{ll}
\text{äußere Symmetrie} &
G_{\mathrm{ext}}=\mathcal P=\R^{1,3}\rtimes SO^+(1,3),\\[0.2em]
\text{Darstellung auf Arena} &
\alpha_{\mathrm{ext}}(a,\Lambda):x\mapsto \Lambda x+a,\\[0.2em]
\text{Darstellung auf Faser} &
\rho_{\mathrm{ext}}(\Lambda):\psi(x)\mapsto S(\Lambda)\psi(\Lambda^{-1}x),\\[0.2em]
\text{innere Symmetrie} &
G_{\mathrm{int}}=\{e\},\text{ (trivial für freies Feld)}.\\[0.2em]
\end{array}
\]
Das Kommutativitätsdiagramm
\[
\begin{array}{ccc}
J^1\mathfrak F(U) & \xrightarrow{D_U} & \mathfrak F(U)\\
\downarrow J^1(\rho_{\mathrm{ext}}(\Lambda)) & & \downarrow \rho_{\mathrm{ext}}(\Lambda)\\
J^1\mathfrak F(\Lambda U) & \xrightarrow{D_{\Lambda U}} & \mathfrak F(\Lambda U)
\end{array}
\]
formalisiert den Satz:
„Die Gestalt der Gleichung ändert sich unter Beobachterwechsel nicht.“
Damit ist der formale Begriff \emph{Invarianz} eine Eigenschaft der
Natürlichkeit des Differentialoperators \(D\).

\bigskip
\emph{9. Fazit.}
\begin{enumerate}
\item Äußere Symmetrien realisieren den Wechsel des Bezugssystems:
\(\Aut_{\mathbf{Geom}}(\mathcal A)\).
\item Innere Symmetrien realisieren den Wechsel der Faserbasis:
\(\Aut_{\mathbf{Bund}}(\mathcal F\mid_{\mathrm{id}_\mathcal A})\).
\item Beide setzen sich zu einer gemeinsamen Gruppoid-Struktur zusammen,
deren Aktion auf der Feldgarbe \(\mathfrak F\) formal die
„Möglichkeit äquivalenter Beschreibungen“ ausdrückt.
\item Invarianz physikalischer Gleichungen ist dann eine Natürlichkeitseigenschaft:
\[
D\circ J^k(\rho(\Lambda,\gamma))
=\rho(\Lambda,\gamma)\circ D.
\]
Dies ist die mathematisch präzise Form der Aussage:
„Das physikalische Gesetz gilt in allen Bezugssystemen gleichermaßen.“
\end{enumerate}


\pagebreak


\subsection{Formale Bedeutung der Invarianz als Natürlichkeitseigenschaft}

\emph{1. Formale Struktur.}
Seien gegeben:
\begin{itemize}
\item eine geometrische Arena $(\mathcal A,g)$ mit Symmetriegruppe $G_{\mathrm{ext}}$,
\item ein Feldbündel $\pi:\mathcal F\to\mathcal A$ mit kombinierter Aktion
      $\rho:G_{\mathrm{ext}}\times G_{\mathrm{int}}\to\Aut_{\mathbf{Bund}}(\mathcal F)$,
\item ein $k$-ter natürlicher Differentialoperator
      \[
      D:J^k\mathcal F\longrightarrow\mathcal E
      \]
      zwischen Bündeln über $\mathcal A$ (im einfachsten Fall $\mathcal E=\mathcal F$).
\end{itemize}
Für $U\subset\mathcal A$ bezeichnet $D_U:\Gamma(U,\mathcal F|_U)\to\Gamma(U,\mathcal E|_U)$
die lokal induzierte Abbildung.  

\textbf{Definition.}
Der Operator $D$ heißt \emph{$G_{\mathrm{tot}}$-kovariant} (oder „invariant in Form“),
wenn für jedes Paar $(\Lambda,\gamma)\in G_{\mathrm{tot}}=G_{\mathrm{ext}}\ltimes G_{\mathrm{int}}$
und alle offenen Mengen $U\subset\mathcal A$ das Diagramm kommutiert:
\[
\begin{array}{ccc}
J^k\mathfrak F(U) & \xrightarrow{\quad D_U\quad} & \mathfrak E(U)\\[0.4em]
\downarrow J^k(\rho(\Lambda,\gamma)) & & \downarrow \rho(\Lambda,\gamma)\\[0.4em]
J^k\mathfrak F(\Lambda U) & \xrightarrow{\quad D_{\Lambda U}\quad} & \mathfrak E(\Lambda U)
\end{array}
\]
Dies ist äquivalent zur Gleichung
\[
D_{\Lambda U}\circ J^k(\rho(\Lambda,\gamma))
= \rho(\Lambda,\gamma)\circ D_U.
\]
Formal ist dies die \emph{Natürlichkeit} von $D$ als Transformation
zwischen Funktoren $\mathfrak F$ und $\mathfrak E$ mit $G_{\mathrm{tot}}$-Aktion:
\[
D:\ J^k\mathfrak F\Rightarrow \mathfrak E
\quad\text{in der Kategorie } [\mathbf{Open}(\mathcal A)^{\mathrm{op}},\mathbf{Set}]^{G_{\mathrm{tot}}}.
\]

\textbf{Interpretation.}
Diese Natürlichkeit bedeutet:
- Die „Gestalt“ des Gesetzes $D(\psi)=0$ bleibt nach einem zulässigen
  Wechsel der Beschreibung unverändert.
- Der Transformationswechsel $(\Lambda,\gamma)$ überführt Lösungen in Lösungen,
  ohne die Funktionsform des Gesetzes zu ändern.

\bigskip
\emph{2. Anwendung: Dirac-Gleichung.}

\textbf{Daten.}
\begin{align*}
\mathcal A &= (\R^{1,3},g),\qquad g=\mathrm{diag}(1,-1,-1,-1),\\
\mathcal F &= M\times\C^4,\qquad
D\psi = i\gamma^\mu\partial_\mu\psi - m\psi.
\end{align*}
\[
\rho_{\mathrm{ext}}(a,\Lambda)\psi(x)
= S(\Lambda)\,\psi(\Lambda^{-1}(x-a)),
\quad
S(\Lambda)=\exp\!\Big(-\tfrac{i}{4}\omega_{\mu\nu}\sigma^{\mu\nu}\Big).
\]

\textbf{Operatorvergleich vor und nach Transformation.}
Wir wollen zeigen:
\[
D'(\psi')(x') = 0 \quad\Longleftrightarrow\quad D(\psi)(x)=0,
\]
d.\,h. die Gleichung behält dieselbe Form in allen Bezugssystemen.

Berechne:
\begin{align*}
(D'\psi')(x')
&= i\gamma^\mu\partial'_\mu\psi'(x') - m\psi'(x')\\
&= i\gamma^\mu\Lambda_\mu{}^\nu\partial_\nu
\big(S(\Lambda)\psi(\Lambda^{-1}(x'-a))\big) - m S(\Lambda)\psi(\Lambda^{-1}(x'-a)).
\end{align*}
Da $S(\Lambda)$ konstant ist (unabhängig von $x$) und
$\gamma^\mu\Lambda_\mu{}^\nu=S(\Lambda)\gamma^\nu S(\Lambda)^{-1}$,
folgt
\begin{align*}
(D'\psi')(x')
&= i S(\Lambda)\gamma^\nu\partial_\nu\psi(\Lambda^{-1}(x'-a)) - mS(\Lambda)\psi(\Lambda^{-1}(x'-a))\\
&= S(\Lambda)\big(i\gamma^\nu\partial_\nu\psi(x)-m\psi(x)\big)\\
&= S(\Lambda)\,(D\psi)(x).
\end{align*}
Damit gilt:
\[
D'(\psi')=S(\Lambda)\,(D\psi)\circ\Lambda^{-1}.
\]
\textbf{Schluss:}
$D\psi=0$ impliziert $D'\psi'=0$,
und die Gleichung bleibt \emph{formgleich} unter Lorentztransformationen.  

Dies ist exakt die Instanziierung der formalen Bedingung
\[
D_{\Lambda U}\circ J^1(\rho_{\mathrm{ext}}(\Lambda))
= \rho_{\mathrm{ext}}(\Lambda)\circ D_U.
\]

\bigskip
\emph{3. Diagrammatische Lesart.}
\[
\begin{array}{ccc}
\psi \in \Gamma(U,\mathcal F) & \mapsto{\ D\ } & D\psi \in \Gamma(U,\mathcal F)\\[0.4em]
\downarrow \rho_{\mathrm{ext}}(\Lambda) & & \downarrow \rho_{\mathrm{ext}}(\Lambda)\\[0.4em]
\psi'=\rho_{\mathrm{ext}}(\Lambda)\psi \in \Gamma(\Lambda U,\mathcal F)
& \mapsto{\ D'\ } &
D'\psi' = \rho_{\mathrm{ext}}(\Lambda)D\psi \in \Gamma(\Lambda U,\mathcal F)
\end{array}
\]
Das Diagramm kommutiert, d.\,h.\ die Transformation „zieht durch den Operator hindurch“:
\[
\rho_{\mathrm{ext}}(\Lambda)\circ D = D\circ \rho_{\mathrm{ext}}(\Lambda).
\]
\emph{Folge:}
Die Gleichung $D\psi=0$ hat dieselbe Form in allen Darstellungen —  
formal: sie ist ein \emph{natürliches Gesetz} im Sinn der Kategorie der Bündel über $\mathcal A$.

\bigskip
\emph{4. Konzeptionelle Deutung.}

Die Natürlichkeitseigenschaft
\[
D\circ J^k(\rho(\Lambda,\gamma)) = \rho(\Lambda,\gamma)\circ D
\]
impliziert:
\begin{enumerate}
\item Die Form der Gesetzesabbildung $D$ hängt nicht von der Wahl des Bezugssystems
      oder der Eichwahl ab.
\item Die Menge der Lösungen $\mathcal E=\ker D$ ist stabil unter allen zulässigen
      Transformationswechseln.
\item Das Gesetz ist somit \emph{objektiv}: seine Struktur ist eine Eigenschaft
      des physikalischen Systems selbst, nicht der Beschreibung.
\end{enumerate}

Im Dirac-Beispiel konkret:
der Operator \(i\gamma^\mu\partial_\mu - m\) enthält nur Tensorobjekte
($\gamma^\mu$ bilden eine Darstellung des Clifford-Algebra-Tensorsystems,
$\partial_\mu$ sind Komponenten eines 4-Vektors);
daher transformiert er kovariant, und die Form der Gleichung bleibt erhalten.
Diese algebraische Invarianz \(\gamma^\mu\Lambda_\mu{}^\nu=S(\Lambda)\gamma^\nu S(\Lambda)^{-1}\)
ist die konkrete Realisierung der Natürlichkeit von \(D\).

\textbf{Zusammenfassung:}
\[
\boxed{
\text{„Natürlichkeit“ $=$ algebraische Forminvarianz der Feldgleichung unter Beschreibungwechseln.}
}
\]

\pagebreak

\subsection{Formalisierung von Beschreibungswechsel und Forminvarianz}

\emph{Ausgangsstruktur.}
Fixiere eine Arena $\mathcal A$ (Objekt in $\mathbf{Geom}$) und ein Feldbündel
$\pi:\mathcal F\to\mathcal A$. Schreibe $\mathfrak F(U)=\Gamma(U,\mathcal F|_U)$
für die Garbe der lokalen Felder. Eine \emph{Beschreibung} ist ein Paar
$(U,\phi)$ mit $U\subset\mathcal A$ offen und $\phi\in\mathfrak F(U)$.

\emph{Kategorie der Beschreibungen.}
Definiere das Gruppoid
\[
\mathsf{Desc}(\mathsf{Sys})
=\Big(\ \mathrm{Obj}=\{(U,\phi)\},\ \mathrm{Mor}((U,\phi),(U',\phi'))=\{(\Lambda,\gamma)\mid \alpha_{\mathrm{ext}}(\Lambda)U=U',\ \rho(\Lambda,\gamma)\phi=\phi'\ \text{auf }U\cap U'\}\ \Big),
\]
wobei $\rho(\Lambda,\gamma):\mathfrak F(U)\to\mathfrak F(\Lambda U)$ der durch äußere
und innere Symmetrien induzierte Feldübersetzer ist.

\emph{Gesetze als natürliche (Differential-)Relationen.}
Ein \emph{Gesetz} ist ein Unterfunktor $\mathcal E\subset \mathfrak F$
(\emph{Lösungsfunktor}) oder äquivalent die Nullstellenmenge eines natürlichen
Differentialoperators $D:J^k\mathfrak F\Rightarrow \mathfrak E$:
\[
\mathcal E(U)=\{\phi\in\mathfrak F(U)\mid D_U(j^k\phi)=0\}.
\]
Die konzeptuelle Rolle von \emph{Gesetz} im abstrakten Setting:
\begin{itemize}
\item \emph{Selektion}: $\mathcal E$ selektiert aus allen Beschreibungen diejenigen, die dem Modell genügen.
\item \emph{Strukturverträglichkeit}: Natürlichkeit von $D$ verankert $\mathcal E$ in den zulässigen Morphismen (Wechseln).
\item \emph{Objektivität}: Forminvarianz (s.u.) macht $\mathcal E$ unabhängig von der Wahl der Beschreibung.
\end{itemize}

\emph{Übersetzungen (Wechsel der Beschreibung).}
Ein Morphismus $(\Lambda,\gamma):(U,\phi)\to(U',\phi')$ besteht aus
\[
U'=\alpha_{\mathrm{ext}}(\Lambda)U,\qquad
\phi'=\rho(\Lambda,\gamma)\phi,
\]
wobei
\[
\rho(\Lambda,\gamma)=\rho_{\mathrm{ext}}(\Lambda)\,\rho_{\mathrm{int}}(\gamma):\ \mathfrak F(U)\to \mathfrak F(U').
\]
Für Jets wird die kanonische Fortsetzung $J^k(\rho(\Lambda,\gamma)):J^k\mathfrak F(U)\to J^k\mathfrak F(U')$ benutzt.

\emph{Forminvarianz als Natürlichkeit.}
Der Operator $D$ heißt \emph{forminvariant} (kovariant) unter den Wechseln, wenn
für alle Morphismen $(\Lambda,\gamma)$ das Diagramm
\[
\begin{array}{ccc}
J^k\mathfrak F(U) & \xrightarrow{\ \ D_U\ \ } & \mathfrak E(U)\\[0.4em]
\downarrow J^k(\rho(\Lambda,\gamma)) & & \downarrow \rho(\Lambda,\gamma)\\[0.4em]
J^k\mathfrak F(U') & \xrightarrow{\ \ D_{U'}\ \ } & \mathfrak E(U')
\end{array}
\]
kommutiert, d.h.
\[
D_{U'}\circ J^k(\rho(\Lambda,\gamma))=\rho(\Lambda,\gamma)\circ D_U.
\]
Dies ist exakt die Natürlichkeit von $D:J^k\mathfrak F\Rightarrow \mathfrak E$
in der funktoriell mit $G_{\mathrm{tot}}$-Aktion versehenen Funktorkategorie.

\emph{Beweisprinzip (Forminvarianz $\Rightarrow$ Lösungserhalt).}
Sei $(U,\phi)$ eine Beschreibung mit erfülltem Gesetz $D_U(j^k\phi)=0$.
Sei $(\Lambda,\gamma):(U,\phi)\to(U',\phi')$ ein Beschreibungswechsel, also $\phi'=\rho(\Lambda,\gamma)\phi$.
Dann gilt:
\[
\underbrace{D_{U'}\big(j^k\phi'\big)}_{\text{Gesetz in neuer Beschreibung}}
= D_{U'}\big(J^k(\rho(\Lambda,\gamma))\,j^k\phi\big)
= \rho(\Lambda,\gamma)\,D_U(j^k\phi)
= \rho(\Lambda,\gamma)\,0
= 0.
\]
Erste Gleichheit: Definition von $D_{U'}$ auf $\phi'$,
zweite: Natürlichkeit,
dritte: Gesetz in alter Beschreibung,
vierte: Funktorialität. Damit ist das Gesetz \emph{formgleich} und lösungserhaltend.

\emph{Lesart im Gruppoid.}
Die Zuordnung
\[
(U,\phi)\ \mapsto\ D_U(j^k\phi)
\]
ist ein natürlicher Transformator $\mathsf{Desc}(\mathsf{Sys})\to \mathbf{Vec}$.
Forminvarianz ist die Aussage, dass diese Zuordnung mit allen Morphismen
$(\Lambda,\gamma)$ kommutiert. Der Lösungsfunktor $\mathcal E=\ker D$ ist dann
ein Untergruppoid von $\mathsf{Desc}(\mathsf{Sys})$.

\emph{Explizites Beispiel: Dirac-Operator.}
Arena $\mathcal A=(\R^{1,3},g)$, Feldgarbe $\mathfrak F(U)=C^\infty(U,\C^4)$,
Operator $D\psi=i\gamma^\mu\partial_\mu\psi - m\psi$.
Ein äußerer Wechsel sei $\Lambda\in SO^+(1,3)$, $(\rho_{\mathrm{int}}=\mathrm{id})$.
Die Übersetzung ist
\[
(\rho(\Lambda)\psi)(x')=S(\Lambda)\,\psi(\Lambda^{-1}x'),\qquad
S(\Lambda)^{-1}\gamma^\mu S(\Lambda)=\Lambda^\mu{}_\nu\,\gamma^\nu.
\]
Berechne:
\[
D_{U'}(\rho(\Lambda)\psi)(x')
= i\gamma^\mu\partial'_\mu\big(S(\Lambda)\psi(\Lambda^{-1}x')\big)-mS(\Lambda)\psi(\Lambda^{-1}x').
\]
Mit $\partial'_\mu=\Lambda_\mu{}^\nu\partial_\nu$ und der Clifford-Kovarianz folgt
\[
= i\,S(\Lambda)\gamma^\nu\partial_\nu\psi(\Lambda^{-1}x')-m\,S(\Lambda)\psi(\Lambda^{-1}x')
= S(\Lambda)\,(D_U\psi)(\Lambda^{-1}x').
\]
Damit
\[
D_{U'}\circ \rho(\Lambda)=\rho(\Lambda)\circ D_U
\quad\Rightarrow\quad
D_U\psi=0\ \Longrightarrow\ D_{U'}(\rho(\Lambda)\psi)=0,
\]
und die Form des Gesetzes ist erhalten.

\emph{Begriffsrelationen (Synopsis).}
\begin{itemize}
\item \emph{Beschreibung}: Objekt $(U,\phi)$ im Gruppoid $\mathsf{Desc}$.
\item \emph{Wechsel/Übersetzung}: Morphismus $(\Lambda,\gamma):(U,\phi)\to(U',\phi')$, realisiert durch $\rho(\Lambda,\gamma)$.
\item \emph{Gesetz}: natürlicher Differentialoperator $D$ (oder dessen Kern $\mathcal E$) als natürlicher Unterfunktor von $\mathfrak F$.
\item \emph{Forminvarianz}: Natürlichkeit von $D$; Kommutativität mit allen Übersetzungen.
\item \emph{Objektivität}: Stabilität des Lösungsgruppoids $\mathcal E$ unter $\mathsf{Desc}$-Morphismen.
\end{itemize}

\emph{Minimalaxiome für Forminvarianz.}
\begin{enumerate}
\item (Garbigkeit) $U\mapsto\mathfrak F(U)$ ist eine Garbe mit Funktorialität $r_{U,V}$.
\item (Äquivariante Aktion) Es existiert $\rho(\cdot):\mathsf{Mor}(\mathsf{Desc})\to \mathrm{Iso}$ mit
$\rho(\mathrm{id})=\mathrm{id}$ und $\rho(\beta\circ\alpha)=\rho(\beta)\circ\rho(\alpha)$.
\item (Natürlichkeit) $D:J^k\mathfrak F\Rightarrow\mathfrak E$ mit
$D_{U'}\circ J^k(\rho)=\rho\circ D_U$ für alle Morphismen $\rho$.
\end{enumerate}
Dann ist $\ker D$ ein unter Morphismen abgeschlossenes Untergruppoid – das \emph{Gesetz} ist formgleich in allen Beschreibungen.

\pagebreak

\subsection{Jets, Natürlichkeit von Differentialoperatoren und Unterfunktoren}

\emph{Ziel.}
Wir präzisieren drei Bausteine des abstrakten Rahmens:
(1) $k$-Jets als kanonische Träger für „bis zur $k$-ten Ableitung“
von Feldern, (2) \emph{Natürlichkeit} lokaler Differentialoperatoren
als Kommutativität mit allen zulässigen Wechseln der Beschreibung,
(3) \emph{Unterfunktor} (Subfunktor) als formale Fassung einer
gleichungsdefinierten Teilgarbe (Lösungsfunktor).

\bigskip
\emph{1. $k$-Jet-Bündel $J^k\mathcal F$.}

\textbf{(a) $k$-Jets von Abschnitten.}
Sei $\pi:\mathcal F\to \mathcal A$ ein glattes Vektorbündel über
einer $n$-dimensionalen glatten Mannigfaltigkeit $\mathcal A$.
Für $x\in\mathcal A$ seien $\phi,\psi\in\Gamma(\mathcal F)$ zwei glatte Abschnitte.
Wir schreiben
\[
j_x^k\phi = j_x^k\psi \quad \Longleftrightarrow \quad
\text{$\phi$ und $\psi$ stimmen in $x$ bis zur $k$-ten Ableitung überein.}
\]
Die Äquivalenzklassen $j_x^k\phi$ heißen \emph{$k$-Jets von $\phi$ in $x$}.
Die Vereinigung aller $k$-Jets
\[
J^k\mathcal F := \bigsqcup_{x\in\mathcal A} J_x^k\mathcal F,
\qquad J_x^k\mathcal F:=\{j_x^k\phi\mid \phi\in\Gamma(\mathcal F)\},
\]
trägt eine kanonische glatte Bündelstruktur
\[
\pi_k:\ J^k\mathcal F \longrightarrow \mathcal A,\qquad j_x^k\phi \mapsto x.
\]

\textbf{(b) Funktorialität (Jet-Prolongation).}
Ein Bündelmorphismus $F:\mathcal F\to\mathcal F'$ über $f:\mathcal A\to\mathcal A'$
induziert den \emph{$k$-Jet-Prolongationsmorphismus}
\[
J^kF:\ J^k\mathcal F \longrightarrow J^k\mathcal F',
\qquad
J^kF\big(j_x^k\phi\big) := j_{f(x)}^k\big(F\circ \phi \circ f^{-1}\big),
\]
sofern $f$ ein Diffeomorphismus (hier in unseren Symmetrieanwendungen: Isomorphie der Arena) ist.
Damit wird $J^k$ zum Funktor auf der Kategorie der Bündelmorphismen.

\textbf{(c) Lokale Beschreibung.}
In lokalen Koordinaten $(x^i)$ und Faserkoordinaten $(u^\alpha)$ mit Partials
$u^\alpha_I = \partial_I u^\alpha$ (Multiindex $I$ mit $|I|\le k$)
werden Punkte von $J^k\mathcal F$ durch
\[
(x^i, u^\alpha, u^\alpha_{i}, u^\alpha_{i j}, \dots, u^\alpha_{I})
\]
parametrisiert. Die Projektion $\pi_k$ vergisst die Ableitungen.

\textbf{(d) Affine/lineare Struktur.}
Für $k\ge 1$ ist $\pi_k:J^k\mathcal F\to \mathcal A$ ein \emph{affines} Bündel
über $J^{k-1}\mathcal F$, modelliert auf
\[
\mathrm{Sym}^k(T^\ast\mathcal A)\otimes \mathcal F,
\]
und $J^0\mathcal F=\mathcal F$. Diese affine Struktur kodiert,
dass die $k$-ten Ableitungen linear aufsetzen.

\bigskip
\emph{2. Lokale Differentialoperatoren als natürliche Transformationen.}

\textbf{(a) Definition (koordinatenfrei).}
Ein \emph{lokaler Differentialoperator der Ordnung $\le k$} von $\mathcal F$
in ein Zielbündel $\mathcal E\to\mathcal A$ ist eine glatte Bündelabbildung
\[
D:\ J^k\mathcal F \longrightarrow \mathcal E,
\]
die sich auf offene Mengen $U\subset\mathcal A$ als Abbildung
\[
D_U:\ \Gamma(U,\mathcal F|_U)\ \longrightarrow\ \Gamma(U,\mathcal E|_U),
\qquad
\phi\ \mapsto\ D\circ j^k\phi
\]
faktorisieren lässt. Hier ist $j^k\phi:U\to J^k\mathcal F|_U$ der \emph{$k$-Jet-Auftrieb}
von $\phi$.

\textbf{(b) Natürlichkeit (Forminvarianz).}
Sei $\rho(\Lambda,\gamma):\mathcal F\to\mathcal F$ ein Bündelautomorphismus
über $\alpha_{\mathrm{ext}}(\Lambda):\mathcal A\to\mathcal A$ (äußerer Wechsel)
mit ggf. vertikaler innerer Wirkung $\gamma$ auf den Fasern.
Dann heißt $D$ \emph{natürlich} (bzw. \emph{forminvariant}) unter $\rho$, wenn
für alle $U\subset\mathcal A$ das Diagramm kommutiert:
\[
\begin{array}{ccc}
J^k\mathfrak F(U) & \xrightarrow{\ \ D_U\ \ } & \mathfrak E(U)\\[0.4em]
\downarrow J^k\big(\rho(\Lambda,\gamma)\big) & & \downarrow \rho(\Lambda,\gamma)\\[0.4em]
J^k\mathfrak F(\Lambda U) & \xrightarrow{\ \ D_{\Lambda U}\ \ } & \mathfrak E(\Lambda U)
\end{array}
\]
d.\,h.
\[
D_{\Lambda U}\circ J^k\big(\rho(\Lambda,\gamma)\big)
=\rho(\Lambda,\gamma)\circ D_U.
\]
Interpretation: „Der Wechsel \emph{zieht durch} den Operator hindurch“,
die Form des Gesetzes bleibt erhalten.

\textbf{(c) Lokale Formulierung.}
In lokalen Koordinaten besitzt $D$ die Form
\[
(D\phi)^\beta(x)=F^\beta\big(x,\,u^\alpha(x),\,\partial_I u^\alpha(x)\big),\quad |I|\le k,
\]
wobei \emph{Natürlichkeit} bedeutet, dass die $F^\beta$ die Tensor-/Darstellungsregeln
unter $(\Lambda,\gamma)$ erfüllen (kein expliziter Koordinaten-/Basisbruch).

\bigskip
\emph{3. Gesetze als Unterfunktoren (Subgarben).}

\textbf{(a) Prägarben- und Garbensicht.}
Der Feldfunktor
\[
\mathfrak F:\ \mathrm{Open}(\mathcal A)^{\mathrm{op}}\longrightarrow \mathbf{Set},\quad
U\mapsto \Gamma(U,\mathcal F|_U)
\]
ist eine Prägarbe (mit Sheaf-Eigenschaft). Ein \emph{Unterfunktor}
\[
\mathcal E\ \subset\ \mathfrak F
\]
ist gegeben durch Teilmengen $\mathcal E(U)\subset\mathfrak F(U)$ für alle
$U$, die \emph{mit den Restriktionen verträglich} sind:
$r_{U,V}(\mathcal E(V))\subset \mathcal E(U)$ für $U\subset V$.
Ist $\mathfrak F$ eine Garbe, so ist $\mathcal E$ eine \emph{Untergarbe},
wenn zusätzlich die Sheaf-Axiome für $\mathcal E$ gelten (lokales Kleben).

\textbf{(b) Gesetz = Nullstellen-Unterfunktor.}
Ein \emph{Gesetz} definieren wir als Nullstellenmenge eines natürlichen
Differentialoperators $D$:
\[
\mathcal E(U)=\big\{\ \phi\in\mathfrak F(U)\ \big|\ D_U(j^k\phi)=0\ \big\}.
\]
Dann ist $\mathcal E$ (i) ein Unterfunktor von $\mathfrak F$ und (ii) $G_{\mathrm{tot}}$-kovariant,
denn Natürlichkeit liefert für $(\Lambda,\gamma)$:
\[
D_{\Lambda U}\big(j^k(\rho(\Lambda,\gamma)\phi)\big)
=\rho(\Lambda,\gamma)\,D_U(j^k\phi).
\]
Insbesondere: $D_U(j^k\phi)=0\ \Rightarrow\ D_{\Lambda U}(j^k\phi')=0$
für $\phi'=\rho(\Lambda,\gamma)\phi$.

\textbf{(c) Morphismus-Sicht (natürliche Monomorphismen).}
Die Inklusionen $\iota_U:\mathcal E(U)\hookrightarrow\mathfrak F(U)$ bilden
einen natürlichen Monomorphismus von Prägarben
\[
\iota:\ \mathcal E \Rightarrow \mathfrak F.
\]
„Unterfunktor“ heißt hier: \emph{natürliche} Teilmengenwahl kompatibel mit allen
Restriktionen; die $G_{\mathrm{tot}}$-Kovarianz macht $\iota$ zu einem
$G_{\mathrm{tot}}$-äquivarianten Monomorphismus.

\bigskip
\emph{4. Zusammenspiel der drei Konzepte im Beweis der Forminvarianz.}

\textbf{(i) Start.}
Wähle eine Beschreibung $(U,\phi)$ und setze voraus, dass das Gesetz auf $U$
erfüllt ist:
\[
\phi\in\mathcal E(U)
\quad\Longleftrightarrow\quad
D_U\big(j^k\phi\big)=0.
\]

\textbf{(ii) Wechsel.}
Nimm einen Morphismus (Beschreibungswechsel)
\[
(\Lambda,\gamma):(U,\phi)\longrightarrow (U',\phi'),\quad
U'=\Lambda U,\ \ \phi'=\rho(\Lambda,\gamma)\phi.
\]
Auf Jet-Ebene wirkt $J^k(\rho(\Lambda,\gamma))$.

\textbf{(iii) Natürlichkeit anwenden.}
Kommutativität liefert:
\[
D_{U'}\big(j^k\phi'\big)
= D_{U'}\circ J^k(\rho(\Lambda,\gamma))\big(j^k\phi\big)
= \rho(\Lambda,\gamma)\circ D_U\big(j^k\phi\big)
= \rho(\Lambda,\gamma)\,0
=0.
\]

\textbf{(iv) Schluss.}
Damit $\phi'\in\mathcal E(U')$.  
Die \emph{Form} der Gleichung (als Gleichheitsrelation im Zielbündel) ist unverändert,
und der Lösungsraum wird durch den Wechsel isomorph auf den Lösungsraum in der
neuen Beschreibung transportiert.  
Formal: $\mathcal E$ ist ein $G_{\mathrm{tot}}$-äquivariant definierter \emph{Unterfunktor}
von $\mathfrak F$, erzeugt als Kern eines \emph{natürlichen} Operators $D$.

\bigskip
\emph{5. Minimaler formaler Kern (in einem Satz).}

$k$-Jets machen „bis zur $k$-ten Ableitung“ funktoriell; \
„Gesetz“ $=$ Kern eines natürlichen $D:J^k\mathfrak F\Rightarrow\mathfrak E$; Unterfunktor $\mathcal E\subset\mathfrak F$ ist dadurch $G_{\mathrm{tot}}$-kovariant und forminvariant.

\pagebreak

\subsection{Systematische und formelle Fassung von Beschreibungswechsel, Pullbacks und Forminvarianz}

\emph{Ziel.}
Wir formulieren den Wechsel von Beschreibungen (Koordinaten/Wahl der Referenzsysteme) und die
Forminvarianz physikalischer Gesetze in einem strikt koordinatenfreien Rahmen. Koordinatenformeln
entstehen dann als Darstellungen der intrinsischen Aussagen.

\bigskip
\emph{1. Kategorien und Grundobjekte.}
\begin{itemize}
\item \emph{Arena-Kategorie} $\mathbf{Geom}$: Objekte sind glatte Mannigfaltigkeiten mit Struktur
(z.\,B. pseudo\-riemann\-sche Metrik $g$); Morphismen sind Struktur-erhaltende Diffeomorphismen.
\item \emph{Bündelkategorie} $\mathbf{Bund}(\mathcal A)$: Objekte sind glatte (Vektor-)Bündel $\pi:\mathcal F\to\mathcal A$;
Morphismen sind Bündelabbildungen über der Identität von $\mathcal A$.
\item \emph{Feldgarbe} $\mathfrak F:\mathrm{Open}(\mathcal A)^{\mathrm{op}}\to\mathbf{Set}$,
$U\mapsto \Gamma(U,\mathcal F|_U)$.
\end{itemize}

\bigskip
\emph{2. Beschreibungswechsel als Pullback.}
\begin{itemize}
\item \emph{Geometrischer Wechsel} ist ein Diffeomorphismus $f:\mathcal A\to\mathcal A$.
\item \emph{Feld-Transformation} (Skalarfeld): $f^*:\ C^\infty(\mathcal A)\to C^\infty(\mathcal A)$, $(f^*\phi)(p):=\phi(f(p))$.
\item \emph{Allgemeines Bündel}: Für ein Bündel $\pi:\mathcal F\to\mathcal A$ entsteht das Pullback-Bündel
$f^*\mathcal F:=\{(p,v)\in \mathcal A\times\mathcal F\mid f(p)=\pi(v)\}\to\mathcal A$ mit kanonischer Projektion.
Abschnitte transformieren via
\[
f^*:\ \Gamma(U,\mathcal F)\longrightarrow \Gamma\big(f^{-1}(U), f^*\mathcal F\big),\qquad
\phi\mapsto f^*\phi.
\]
\item \emph{Isometrien}: Ist $g$ eine Pseudo\-metrik, so induziert $f$ die Metrik $f^*g$; bei Isometrien gilt $f^*g=g$.
\end{itemize}

\bigskip
\emph{3. Natürliche Differentialoperatoren.}
\begin{itemize}
\item \emph{Jet-Funktor}: Für ein Bündel $\mathcal F\to\mathcal A$ sei $J^k\mathcal F\to\mathcal A$ das $k$-Jet-Bündel.
Jeder Bündelmorphismus $F:\mathcal F\to\mathcal F'$ über $f:\mathcal A\to\mathcal A'$ induziert
$J^kF:J^k\mathcal F\to J^k\mathcal F'$ (Funktorialität).
\item \emph{Lokaler Operator Ordnung $\le k$}:
eine glatte Bündelabbildung $D:J^k\mathcal F\to\mathcal E$ (Zielbündel $\mathcal E\to\mathcal A$), die lokal
$D_U:\Gamma(U,\mathcal F)\to\Gamma(U,\mathcal E)$ durch $D_U(\phi)=D\circ j^k\phi$ faktorisieren lässt.
\item \emph{Natürlichkeit (Forminvarianz)}:
Für jeden Diffeomorphismus $f$ ist $D$ \emph{natürlich}, falls
\[
\boxed{ \quad D\_{f(U)}\circ J^k(f^*) \;=\; f^*\circ D_U \quad }
\]
als Abbildungen $J^k\Gamma(U,\mathcal F)\to \Gamma(f(U),\mathcal E)$ gilt. Diagrammatisch:
\[
\begin{array}{ccc}
J^k\mathfrak F(U) & \xrightarrow{\ \ D_U\ \ } & \mathfrak E(U)\\[0.2em]
\downarrow J^k(f^*) & & \downarrow f^*\\[0.2em]
J^k\mathfrak F(f(U)) & \xrightarrow{\ \ D_{f(U)}\ \ } & \mathfrak E(f(U))
\end{array}
\]
\end{itemize}

\bigskip
\emph{4. Gesetze als Unterfunktoren.}
\begin{itemize}
\item \emph{Gesetz} (auf $\mathcal A$) ist der Nullstellen-Unterfunktor eines natürlichen Operators $D$:
\[
\mathcal E(U):=\{\ \phi\in\mathfrak F(U)\mid D_U(j^k\phi)=0\ \}\ \subset\ \mathfrak F(U).
\]
\item \emph{Kovarianz des Gesetzes}: Aus Natürlichkeit folgt
\[
f^*(\mathcal E(U))=\mathcal E(f(U)) \quad\text{(Lösungen werden auf Lösungen abgebildet).}
\]
\end{itemize}

\bigskip
\emph{5. Koordinaten als Darstellungen.}
\begin{itemize}
\item Karten $x:U\to\R^n$ und $x':U'\to\R^n$ sind Darstellungen. Der Kartenwechsel ist
$\Phi:=x'\circ x^{-1}:\ x(U\cap U')\to x'(U\cap U')$.
\item Für einen geometrischen Wechsel $f:\mathcal A\to\mathcal A$ wähle $x' := x\circ f^{-1}$.
Dann ist die \emph{Darstellung} des Pullbacks $f^*\phi$ genau die übliche Formel
\[
\phi'(x')=(f^*\phi)(x')=\phi\big(x\big)=\phi\big(\Phi^{-1}(x')\big).
\]
Differentialoperatoren in Koordinaten entstehen durch Kettenregel aus der intrinsischen
Natürlichkeit (kein „Variablen-Mixen“ nötig).
\end{itemize}

\bigskip
\emph{6. Beispiel A: Klein–Gordon operatoriell (koordinatenfrei).}
\begin{itemize}
\item \emph{Operator}: $\square_g:=\mathrm{div}_g\circ\nabla^g = g^{ab}\nabla_a\nabla_b$ auf $C^\infty(\mathcal A)$.
\item \emph{Naturtreue}: Für jeden Diffeomorphismus $f$ gilt
\[
\boxed{\quad f^*\big(\square_g\,\phi\big)\ =\ \square_{f^*g}\,\big(f^*\phi\big). \quad}
\]
\item \emph{Isometrie-Fall}: $f^*g=g \ \Rightarrow\  f^*\circ(\square_g+m^2)=(\square_g+m^2)\circ f^*$.
Damit ist $\mathcal D_g:=\square_g+m^2$ natürlich und das Gesetz
$\mathcal E_g:=\ker\mathcal D_g$ kovariant.
\end{itemize}

\bigskip
\emph{7. Beispiel B: Dirac operatoriell (Strichskizze).}
\begin{itemize}
\item \emph{Spin-Hauptbündel} $P_{\mathrm{Spin}}\to(\mathcal A,g)$, zugehöriges Spinorbündel
$\Sigma\to\mathcal A$, Levi-Civita-Spinverbindung $\nabla^\Sigma$.
\item \emph{Clifford-Multiplikation} $c:T^*\mathcal A\otimes\Sigma\to\Sigma$,
Dirac-Operator $\slashed{D}:=c\circ\nabla^\Sigma:\Gamma(\Sigma)\to\Gamma(\Sigma)$.
\item \emph{Natürliche Kovarianz}: Für Isometrie $f$ hebt sich $f$ zu einem
$f^\Sigma:\Sigma\to\Sigma$ an (über die Spinstruktur) und es gilt
\[
\boxed{\quad f^{\Sigma\,\*}\big(\slashed{D}\,\psi\big)\ =\ \slashed{D}\,\big(f^{\Sigma\,\*}\psi\big). \quad}
\]
Mit Masse $m$: $(\slashed{D}-m)$ ist natürlich, das Dirac-Gesetz $\ker(\slashed{D}-m)$ ist kovariant.
\end{itemize}

\bigskip
\emph{8. Formeller Nachweis der Forminvarianz (Schema).}
Sei $D:J^k\mathcal F\to\mathcal E$ natürlich und $(U,\phi)$ mit $D_U(j^k\phi)=0$. Für Diffeo $f$:
\[
D_{f(U)}\big(j^k(f^*\phi)\big)
= D_{f(U)}\big(J^k(f^*)\,j^k\phi\big)
= f^*\big(D_U(j^k\phi)\big)
= f^*(0)=0.
\]
Damit $f^*\phi\in\mathcal E(f(U))$. \emph{Gesetz bleibt in Form gleich.}

\bigskip
\emph{9. Koordinaten-Korollar (Darstellung der obigen Aussagen).}
Wähle Karten $x$ und $x':=x\circ f^{-1}$. Für ein Skalarfeld gilt in Darstellung:
\[
\partial'_\mu = \frac{\partial x^\nu}{\partial x'^\mu}\,\partial_\nu,\qquad
\phi'(x')=\phi\big(\Phi^{-1}(x')\big).
\]
Ist $f$ Lorentz (Isometrie), folgt unmittelbar
\[
(\square'+m^2)\phi' \ =\ \big((\square+m^2)\phi\big)\circ f^{-1},
\]
und analog (über Spin-Lift) für den Dirac-Operator:
\[
(i\gamma^\mu\partial'_\mu - m)\psi' \ =\ S(\Lambda)\,(i\gamma^\mu\partial_\mu - m)\psi\circ f^{-1}.
\]

\bigskip
\emph{10. Zusammenfassung (Minimalaxiome).}
\begin{enumerate}
\item (\emph{Garbigkeit}) $U\mapsto \mathfrak F(U)$ ist eine Garbe.
\item (\emph{Pullback-Funktorialität}) Diffeomorphismen $f$ induzieren $f^*:\mathfrak F\Rightarrow \mathfrak F$.
\item (\emph{Natürlichkeit}) $D:J^k\mathfrak F\Rightarrow \mathfrak E$ erfüllt $D\circ J^k(f^*)=f^*\circ D$.
\end{enumerate}
Dann ist das \emph{Gesetz} $\mathcal E=\ker D$ ein $f^*$-invarianter Unterfunktor; in Koordinaten
erscheint dies als \emph{Forminvarianz} der Gleichung unter zulässigen Beschreibungswechseln.


\pagebreak

\subsection{Ein formeller Rahmen für Beschreibungswechsel und Variablenreferenz}

\emph{Motivation.}
In physikalischen Formeln wie $\phi(x)$ wird $x$ oft stillschweigend als
„Raumzeitpunkt“ aufgefasst. Formal ist dies jedoch eine \emph{Koordinatenwahl}
oder \emph{Darstellungsvorschrift} eines Punktes $p\in M$ durch ein Symbol
$x$. Wir entwickeln einen Formalismus, der diese Unterscheidung explizit macht.

\bigskip
\emph{1. Geometrischer Grundraum.}
\begin{itemize}
\item Sei $M$ eine glatte Mannigfaltigkeit („Raumzeit“).
\item Punkte $p\in M$ sind abstrakte Elemente, \emph{nicht} identisch mit Koordinaten.
\end{itemize}

\bigskip
\emph{2. Koordinaten als Repräsentationsvorschriften.}
\begin{itemize}
\item Eine \emph{Darstellungsvorschrift} (Koordinatenkarte) ist eine glatte Abbildung
\[
x:U\subset M \longrightarrow \mathbb{R}^n,
\]
die jedem Punkt $p\in U$ einen Zahlenvektor $x(p)$ zuweist.
\item $x$ ist kein intrinsischer Bestandteil des Systems, sondern ein
\emph{Bezeichnungsmechanismus}.  
Man schreibt \(\phi(x)\) nur dann, wenn zuvor
\[
\phi(x) := \phi\circ x^{-1},
\]
also die Darstellung der Funktion $\phi:M\to V$ in der Karte $x$ definiert wurde.
\end{itemize}

\bigskip
\emph{3. Beschreibungen als Paare.}
\begin{itemize}
\item Eine \emph{Beschreibung} ist ein Paar
\[
\mathbf{B} = (U,x),
\]
bestehend aus einem offenen Bereich $U\subset M$ und einer Darstellungsvorschrift $x:U\to\R^n$.
\item Eine Feldbeschreibung in dieser Karte ist
\[
\phi^{[\mathbf B]} := \phi\circ x^{-1} : x(U)\to V,
\]
die in Komponenten durch die Variable $x\in\R^n$ beschrieben wird.
\end{itemize}

\bigskip
\emph{4. Beschreibungswechsel als Morphismus.}
\begin{itemize}
\item Zwei Beschreibungen $\mathbf B=(U,x)$ und $\mathbf B'=(U',x')$
seien gegeben mit $U\cap U'\ne\emptyset$.
\item Der \emph{Koordinatenwechsel} (oder „Übersetzungsmechanismus“) zwischen beiden ist die Abbildung
\[
\Phi_{\mathbf{B}\to\mathbf{B}'} := x'\circ x^{-1}:\ x(U\cap U')\longrightarrow x'(U\cap U').
\]
\item Der \emph{geometrische Beschreibungswechsel} (z.\,B. Poincaré-Transformation) sei ein Diffeomorphismus
\[
f:M\to M.
\]
Er wirkt auf Beschreibungen durch
\[
f\cdot(U,x) := \big(f(U),\ x\circ f^{-1}\big),
\]
und auf Felder durch Pullback
\[
(f^*\phi)(p) := \phi(f(p)).
\]
\end{itemize}

\bigskip
\emph{5. Evaluationsregel.}
Um den Wert einer Feldbeschreibung in einer bestimmten Variablen zu berechnen,
muss explizit angegeben werden, in welchem \emph{Darstellungssystem} die Auswertung erfolgt:
\[
\mathrm{ev}_{\mathbf B}(\phi,p) := \phi^{[\mathbf B]}(x(p)) = \phi(p),
\]
wobei $x$ die aktuell gültige Darstellung ist.  
Das Symbol $x$ ist somit kein geometrisches Argument, sondern die Repräsentation des Punktes $p$
unter der aktiven Beschreibung $\mathbf B$.

\bigskip
\emph{6. Umweisung der Variablen.}
Ein Beschreibungswechsel $\mathbf B\to\mathbf B'$ ändert die Zuordnung von Symbolen zu Punkten:
\[
x' = x\circ f^{-1},\qquad
p' = f(p).
\]
Damit transformiert die Darstellung des Feldes nach der Regel
\[
\phi^{[\mathbf B']}(x') 
= (\phi')\circ (x')^{-1}(x')
= (\phi\circ f)\circ f^{-1}\circ x^{-1}(x')
= \phi^{[\mathbf B]}(x),
\]
mit $x=x'\circ f$ als Komposition der Vorschriften.

In Worten:
\emph{Der gleiche geometrische Punkt wird durch unterschiedliche Variablen dargestellt,
aber die zusammengesetzte Auswertung führt zum gleichen Funktionswert.}

\bigskip
\emph{7. Koordinatenfreie Form der Transformation.}
Sei $f:M\to M$ Diffeomorphismus, $x' = x\circ f^{-1}$.
Dann gilt:
\[
\boxed{
\phi^{[\mathbf B']}(x') = (\phi^{[\mathbf B]}\circ \Phi_{\mathbf B'\to\mathbf B})(x'),
\qquad
\Phi_{\mathbf B'\to\mathbf B} = x\circ f^{-1}\circ x'^{-1}.
}
\]
Dies ist die exakte formale Grundlage der in der Physik üblichen Schreibweise
\(\phi'(x') = \phi(x)\), die kurz für
\[
\phi'(x') = (\phi\circ f)(p') = \phi(p),
\]
steht, wobei \(p'=f^{-1}(p)\) und die jeweilige Darstellung implizit ist.

\bigskip
\emph{8. Operatorische Transformation.}
Ein Differentialoperator $D$ auf $C^\infty(M,V)$ (z.\,B. $\Box_g+m^2$) ist \emph{natürlich}, wenn
\[
f^*\circ D = D\circ f^*.
\]
Im Darstellungsformalismus entspricht dies:
\[
D^{[\mathbf B']}\big(\phi^{[\mathbf B']}\big)
= D^{[\mathbf B]}\big(\phi^{[\mathbf B]}\big)\circ \Phi_{\mathbf B'\to\mathbf B}.
\]
Damit:
\[
D^{[\mathbf B']}\big(\phi^{[\mathbf B']}\big)(x') = 0
\quad\Longleftrightarrow\quad
D^{[\mathbf B]}\big(\phi^{[\mathbf B]}\big)(x)=0.
\]
Der \emph{Gesetzesausdruck bleibt formgleich}, wenn die Übersetzungsregeln für Variablen
und Funktionsauswertungen explizit berücksichtigt werden.

\bigskip
\emph{9. Interpretation.}
\begin{itemize}
\item $x$ ist nicht „der Punkt selbst“, sondern eine \emph{Darstellungsvorschrift}.
\item $\phi$ ist eine geometrische Funktion, deren \emph{Darstellung} $\phi^{[\mathbf B]}$
erst durch Wahl von $x$ entsteht.
\item Der Übergang $x\to x'$ ist eine Änderung der symbolischen Referenz, nicht der
physikalischen Größe.
\item Die Forminvarianz folgt aus der Funktorialität der Pullbacks: Die Gleichungen
werden nur dann identisch aussehen, wenn alle symbolischen Umweisungen korrekt
in den Evaluationsketten mitgeführt werden.
\end{itemize}

\bigskip
\emph{10. Schema (kommutatives Diagramm).}

\[
\begin{array}{ccccc}
p & \mapsto{f} & f(p) & & \\
\downarrow x & & \downarrow x' & & \\
x(p) & \mapsto{\Phi_{\mathbf B\to\mathbf B'}} & x'(f(p)) & \text{in Variablenebene} &
\end{array}
\]
und
\[
\phi^{[\mathbf B]}(x(p)) = \phi(p) = \phi'(f(p)) = \phi^{[\mathbf B']}(x'(f(p))).
\]

\textbf{Schlussfolgerung:}
\[
\boxed{
\text{In jedem Ausdruck $\phi(x)$ ist $x$ eine gewählte Repräsentationsvorschrift; 
ein Beschreibungswechsel ersetzt simultan die Zuordnung $x$ und den Feldfunktor.}
}
\]
Dies liefert eine formelle Grundlage, um „naive Variablen“ stets als 
\emph{parametrisierte Bezeichnungsfunktionen} zu behandeln, die explizit
in die Transformation mit einbezogen werden.



\subsection{Explizite Fassung des Beschreibungswechsels auf $(\R^{1,3},g)$ unter Poincaré-Transformationen}

\emph{Ziel.}
Wir wenden den zuvor entwickelten formellen Rahmen konkret auf die spezielle Relativitätstheorie an:
Raumzeit $(M,g)=(\R^{1,3},\eta)$ mit Minkowski-Metrik $\eta=\mathrm{diag}(1,-1,-1,-1)$
und Symmetriegruppe der Poincaré-Transformationen.

\bigskip
\emph{1. Geometrische und algebraische Ausgangsdaten.}
\begin{itemize}
\item \textbf{Arena:} $M=\R^{1,3}$, Punkte $p=(t,\mathbf x)\in\R^4$.
\item \textbf{Metrik:} $\eta(v,w)=v^0w^0-\mathbf v\cdot\mathbf w$.
\item \textbf{Symmetriegruppe:} 
\[
\mathcal P := \R^{1,3}\rtimes SO^+(1,3),
\quad
(a,\Lambda)\cdot (a',\Lambda')=(a+\Lambda a',\Lambda\Lambda').
\]
\item \textbf{Wirkung auf Punkte:}
\[
f_{(a,\Lambda)}:M\longrightarrow M,
\qquad
f_{(a,\Lambda)}(p)=\Lambda p + a.
\]
\end{itemize}

\bigskip
\emph{2. Beschreibungen und Koordinatenvorschriften.}
\begin{itemize}
\item Eine \emph{Beschreibung} ist ein Paar $\mathbf B=(U,x)$ mit $U\subset M$ offen und
$x:U\to\R^4$ einer glatten Vorschrift, die jedem Punkt $p\in U$ Koordinaten $x(p)$ zuordnet.
\item Standardbeschreibung: $\mathbf B_0=(M,\mathrm{id})$ mit $x(p)=p$ (Inertialsystem $S$).
\item Transformierte Beschreibung (neues Bezugssystem):
\[
\mathbf B_{(a,\Lambda)} := f_{(a,\Lambda)}\cdot\mathbf B_0
= \big(M,\ x_{(a,\Lambda)}:=x\circ f_{(a,\Lambda)}^{-1}\big),
\]
d.h.
\[
x_{(a,\Lambda)}(p) = \Lambda^{-1}(p - a).
\]
Dies sind die „Koordinaten im System $S'$“, das relativ zu $S$ um $(a,\Lambda)$ transformiert ist.
\end{itemize}

\bigskip
\emph{3. Felder und ihre Darstellungen.}
\begin{itemize}
\item Sei $\phi:M\to V$ ein (z.\,B. komplexes skalares) Feld.
\item In der Beschreibung $\mathbf B=(U,x)$ lautet die Darstellung:
\[
\phi^{[\mathbf B]} := \phi\circ x^{-1}: x(U)\subset\R^4 \to V.
\]
\item Speziell für $\mathbf B_0=(M,\mathrm{id})$:
\[
\phi^{[\mathbf B_0]}(x) = \phi(x),
\]
die „Standarddarstellung“.
\end{itemize}

\bigskip
\emph{4. Beschreibungswechsel durch $(a,\Lambda)$.}
\begin{itemize}
\item \textbf{Geometrischer Effekt:}
$f_{(a,\Lambda)}:p\mapsto p'=\Lambda p + a$.
\item \textbf{Feld-Pullback:}
\[
(f_{(a,\Lambda)}^*\phi)(p) = \phi(f_{(a,\Lambda)}(p))
= \phi(\Lambda p + a).
\]
Dies ist das „gleiche physikalische Feld“, gesehen vom neuen Bezugssystem.
\item \textbf{Darstellung in der neuen Beschreibung} $\mathbf B_{(a,\Lambda)}=(M,x_{(a,\Lambda)})$:
\[
\phi^{[\mathbf B_{(a,\Lambda)}]}(x')
:= (f_{(a,\Lambda)}^*\phi)\circ (x_{(a,\Lambda)})^{-1}(x')
= \phi\circ f_{(a,\Lambda)}\circ f_{(a,\Lambda)}(x')
= \phi(\Lambda x' + a).
\]
\item \textbf{Wechselabbildung zwischen Koordinaten:}
\[
\Phi_{\mathbf B_0\to \mathbf B_{(a,\Lambda)}} = x_{(a,\Lambda)}\circ x^{-1} = \Lambda^{-1}(x - a).
\]
\end{itemize}

\bigskip
\emph{5. Zusammenhänge der Darstellungen.}
Die beiden Beschreibungen stehen im formalen Verhältnis:
\[
\boxed{
\phi^{[\mathbf B_{(a,\Lambda)}]}(x')
= \phi^{[\mathbf B_0]}\big(\Lambda x' + a\big).
}
\]
Diese Relation ist die koordinatenfreie Form der physikalisch üblichen Schreibweise:
\[
\phi'(x') = \phi(x),\qquad x = \Lambda^{-1}(x'-a).
\]

\bigskip
\emph{6. Ableitungen und Operatoren.}
\begin{itemize}
\item Ableitungsoperatoren transformieren nach der Kettenregel:
\[
\partial'_\mu = \frac{\partial x^\nu}{\partial x'^\mu}\,\partial_\nu
= (\Lambda^{-1})^\nu{}_\mu\,\partial_\nu.
\]
\item Damit wirkt der D’Alembert-Operator
\[
\Box' = \eta^{\mu\nu}\partial'_\mu\partial'_\nu
= \eta^{\mu\nu}(\Lambda^{-1})^\alpha{}_\mu(\Lambda^{-1})^\beta{}_\nu \partial_\alpha\partial_\beta
= \eta^{\alpha\beta}\partial_\alpha\partial_\beta = \Box,
\]
weil $\Lambda$ Lorentz ist ($\Lambda^T\eta\Lambda=\eta$).
\item Der Operator $(\Box+m^2)$ ist daher \emph{natürlich}:
\[
D_{\mathbf B_{(a,\Lambda)}}\circ J^2(f_{(a,\Lambda)}^*)
= f_{(a,\Lambda)}^*\circ D_{\mathbf B_0}.
\]
\end{itemize}

\bigskip
\emph{7. Beispiel 1 – Translation.}
\begin{itemize}
\item Transformation: $(a,\Lambda)=(a,\mathrm{id})$, also $p'=p+a$.
\item Feldabbildung:
\[
\phi'(x')=\phi(x' + a).
\]
\item Ableitungen:
\(\partial'_\mu=\partial_\mu\), also $\Box'=\Box$.
\item Gesetz:
\[
(\Box'+m^2)\phi'(x')=(\Box+m^2)\phi(x'+a)=0
\quad\text{falls}\quad
(\Box+m^2)\phi=0.
\]
\item \textbf{Interpretation:}
Reine Translation verändert nur den Bezugspunkt, nicht die Struktur der Gleichung.
\end{itemize}

\bigskip
\emph{8. Beispiel 2 – Lorentz-Boost.}
\begin{itemize}
\item Boost in $z$-Richtung mit Geschwindigkeit $v$, Lorentzmatrix
\[
\Lambda(v)=
\begin{pmatrix}
\gamma & 0 & 0 & -\gamma v\\
0 & 1 & 0 & 0\\
0 & 0 & 1 & 0\\
-\gamma v & 0 & 0 & \gamma
\end{pmatrix},
\quad \gamma=(1-v^2)^{-1/2}.
\]
\item Transformation des Feldes:
\[
\phi'(t',x',y',z')=\phi(\gamma(t'+vz'),\,x',\,y',\,\gamma(z'+vt')).
\]
\item Kettenregel:
\[
\partial'_{t'}=\gamma(\partial_t+v\partial_z),\quad
\partial'_{z'}=\gamma(\partial_z+v\partial_t),
\quad
\partial'_{x'}=\partial_x,\ \partial'_{y'}=\partial_y.
\]
Damit:
\[
\Box'=\partial'^2_{t'}-\partial'^2_{x'}-\partial'^2_{y'}-\partial'^2_{z'}
=\partial_t^2-\partial_x^2-\partial_y^2-\partial_z^2=\Box.
\]
\item \textbf{Folge:}
\[
(\Box'+m^2)\phi'=(\Box+m^2)\phi\circ f_{(a,\Lambda)}=0.
\]
\end{itemize}

\bigskip
\emph{9. Beispiel 3 – Spinorische Erweiterung.}
\begin{itemize}
\item Spinor-Feld $\psi:M\to\C^4$.
\item Transformation:
\[
\psi'(x') = S(\Lambda)\,\psi(\Lambda^{-1}(x'-a)),
\quad
S(\Lambda)^{-1}\gamma^\mu S(\Lambda)=\Lambda^\mu{}_\nu\gamma^\nu.
\]
\item Differentialoperator:
\[
D = i\gamma^\mu\partial_\mu - m.
\]
\item Kettenregel und Clifford-Kovarianz ergeben:
\[
D'\psi'(x') = S(\Lambda)\,(D\psi)(\Lambda^{-1}(x'-a)),
\]
und somit:
\[
(D'\psi')(x')=0\quad\Longleftrightarrow\quad(D\psi)(x)=0.
\]
\item \textbf{Forminvarianz:}
Der Dirac-Operator ist natürlich unter der Spin-Lift-Aktion der Lorentzgruppe.
\end{itemize}

\bigskip
\emph{10. Fazit.}
In diesem Formalismus werden alle Ausdrücke wie $\phi(x)$ oder $\psi(x)$
erst sinnvoll, nachdem ein Paar $(U,x)$, also eine Beschreibung, fixiert wurde.
Ein Beschreibungswechsel $(a,\Lambda)$:
\begin{itemize}
\item verändert simultan die Zuordnung der Variablen $x\mapsto x'=\Lambda^{-1}(x-a)$,
\item ändert die Funktionsdarstellung durch Pullback $\phi\mapsto f_{(a,\Lambda)}^*\phi$,
\item erhält den geometrischen Operator $(\Box+m^2)$ bzw.\ $(i\slashed{\partial}-m)$ invariant,
\item und garantiert durch Natürlichkeit:
\[
D_{(a,\Lambda)}\big(f_{(a,\Lambda)}^*\phi\big)
= f_{(a,\Lambda)}^*(D\phi),
\]
was sich in Koordinaten als
\[
(D'\phi')(x') = (D\phi)(x)
\]
darstellt — die präzise Bedeutung der üblichen Forminvarianz.
\end{itemize}

\[
\boxed{
\text{Poincaré-Beschreibungswechsel = simultane Änderung von (i) Variablenvorschrift, (ii) Funktionsdarstellung, und (iii) Pullback des Operators.}
}
\]


\pagebreak



\subsection{Zwei explizite Beispiele im vollen Formalismus}
Wir rechnen die Natürlichkeit (Forminvarianz) \emph{explizit} durch für
\begin{enumerate}
\item das \textbf{Klein–Gordon-Feld} (2.\ Ordnung: $k=2$), und
\item das \textbf{Dirac-Feld} (1.\ Ordnung: $k=1$).
\end{enumerate}
Arena jeweils $\mathcal A=(M,g)$ mit $M=\R^{1,3}$, $g=\mathrm{diag}(1,-1,-1,-1)$,
äußere Symmetrie $G_{\mathrm{ext}}=\mathcal P=\R^{1,3}\rtimes SO^+(1,3)$.
Wir notieren den \emph{Beschreibungwechsel}
\[
x'=\alpha_{\mathrm{ext}}(a,\Lambda)(x)=\Lambda x + a,
\qquad (\Lambda\in SO^+(1,3),\ a\in\R^{1,3}).
\]

\bigskip
\emph{Vorbereitung: Jet-Prolongation und Übersetzer.}

Sei $\pi:\mathcal F\to M$ ein (triviales) Feldbündel mit Faser $F_0$ und Feldgarbe
$\mathfrak F(U)=\Gamma(U,\mathcal F|_U)$.

\textbf{(1) Jet-Prolongation.} Für $k\in\N$ ist $J^k\mathcal F\to M$ das $k$-Jet-Bündel.
Für $\phi\in\mathfrak F(U)$ ist $j^k\phi:U\to J^k\mathcal F|_U$ der $k$-Jet-Auftrieb.
Ein Bündelautomorphismus $\rho(\Lambda):\mathcal F\to\mathcal F$ über $\alpha_{\mathrm{ext}}(\Lambda)$
induziert $J^k(\rho(\Lambda)) : J^k\mathcal F\to J^k\mathcal F$ mit
\[
J^k(\rho(\Lambda))(j_x^k\phi)=j_{\Lambda x+a}^k\!\big(\rho(\Lambda)\phi\big).
\]

\textbf{(2) Feldübersetzer (äußere Aktion).} 
\[
\rho_{\mathrm{ext}}(a,\Lambda):\ \mathfrak F(U)\longrightarrow \mathfrak F(\Lambda U + a),
\qquad
(\rho_{\mathrm{ext}}(a,\Lambda)\phi)(x') := \Phi(\Lambda)\,\phi(\Lambda^{-1}(x'-a)).
\]
Dabei ist $\Phi(\Lambda)$ die (evtl. triviale) Darstellungswirkung auf der Faser.
Für Skalare: $\Phi(\Lambda)=\mathrm{id}$;
für Dirac-Spinoren: $\Phi(\Lambda)=S(\Lambda)$ mit $S(\Lambda)^{-1}\gamma^\mu S(\Lambda)=\Lambda^\mu{}_\nu\gamma^\nu$.

\textbf{(3) Natürlichkeit (Forminvarianz).}
Ein lokaler Differentialoperator $D:J^k\mathcal F\to\mathcal E$ heißt \emph{natürlich}
unter $(a,\Lambda)$, wenn
\[
D_{\Lambda U + a}\circ J^k(\rho_{\mathrm{ext}}(a,\Lambda))
=\rho_{\mathrm{ext}}(a,\Lambda)\circ D_U.
\]
Dies ist die Kommutativität des Diagramms
\[
\begin{array}{ccc}
J^k\mathfrak F(U) & \xrightarrow{D_U} & \mathfrak E(U)\\
\downarrow J^k(\rho_{\mathrm{ext}}(a,\Lambda)) & & \downarrow \rho_{\mathrm{ext}}(a,\Lambda)\\
J^k\mathfrak F(\Lambda U + a) & \xrightarrow{D_{\Lambda U+a}} & \mathfrak E(\Lambda U+a)
\end{array}
\]
und bedeutet: „Die Form des Gesetzes bleibt unter dem Wechsel erhalten.“

\bigskip
\subsubsection*{A. Beispiel (2.\ Ordnung): Klein–Gordon}

\emph{Feld, Operator, Gesetz.}
\[
\mathcal F=M\times\C,\quad \mathfrak F(U)=C^\infty(U,\C),
\qquad
D\phi := (\Box + m^2)\phi,
\quad \Box = g^{\mu\nu}\partial_\mu\partial_\nu.
\]
Dies ist ein Operator $D:J^2\mathcal F\to \mathcal F$ (Ordnung $k=2$).
Gesetz: $\mathcal E(U)=\{\phi\in\mathfrak F(U)\mid D_U(j^2\phi)=0\}$.

\emph{Übersetzer (äußere Symmetrie).} Für Skalare $\Phi(\Lambda)=\mathrm{id}$,
also
\[
(\rho_{\mathrm{ext}}(a,\Lambda)\phi)(x')=\phi(\Lambda^{-1}(x'-a)).
\]

\emph{Expliziter Nachweis der Natürlichkeit.}
Wir zeigen
\[
D_{\Lambda U+a}\big(\rho_{\mathrm{ext}}(a,\Lambda)\phi\big)
=\rho_{\mathrm{ext}}(a,\Lambda)\big(D_U\phi\big),
\quad\text{d.\,h.}\quad
(\Box_{x'}+m^2)\phi'(x') = (\Box_x+m^2)\phi(x)\ \circ\ \Lambda^{-1}(\cdot-a),
\]
mit $\phi'=\rho_{\mathrm{ext}}(a,\Lambda)\phi$ und $x=\Lambda^{-1}(x'-a)$.

Kettenregel (Translationsanteil ist trivial, Lorentzanteil linear):
\[
\partial'_\mu = \frac{\partial x^\nu}{\partial x'^\mu}\, \partial_\nu
= (\Lambda^{-1})^\nu{}_\mu\,\partial_\nu.
\]
Dann
\[
\partial'_\mu\partial'_\nu \phi'(x')
= (\Lambda^{-1})^\alpha{}_\mu(\Lambda^{-1})^\beta{}_\nu
\,\partial_\alpha\partial_\beta\ \phi(x).
\]
Kontrahiere mit $g^{\mu\nu}$ und nutze Lorentz-Invarianz $g^{\mu\nu}(\Lambda^{-1})^\alpha{}_\mu(\Lambda^{-1})^\beta{}_\nu
= g^{\alpha\beta}$:
\[
\Box_{x'}\phi'(x')
= g^{\mu\nu}\partial'_\mu\partial'_\nu \phi'(x')
= g^{\alpha\beta}\partial_\alpha\partial_\beta \phi(x)
= (\Box_x\phi)(x).
\]
Damit
\[
(\Box_{x'}+m^2)\phi'(x') = \big((\Box_x+m^2)\phi\big)(x).
\]
Dies ist genau die Gleichung
\[
D_{\Lambda U+a}\circ J^2(\rho_{\mathrm{ext}}(a,\Lambda))
=\rho_{\mathrm{ext}}(a,\Lambda)\circ D_U.
\]
\emph{Folge (Lösungserhalt):}
Ist $D_U(j^2\phi)=0$ auf $U$, so gilt
$D_{\Lambda U+a}\big(j^2(\rho_{\mathrm{ext}}(a,\Lambda)\phi)\big)=0$ auf $\Lambda U+a$.
Die Gesetzesform bleibt unverändert.

\bigskip
\subsubsection*{B. Beispiel (1.\ Ordnung): Dirac}

\emph{Feld, Operator, Gesetz.}
\[
\mathcal F=M\times\C^4,\quad \mathfrak F(U)=C^\infty(U,\C^4),
\qquad
D\psi := i\gamma^\mu\partial_\mu\psi - m\psi,
\]
$D:J^1\mathcal F\to \mathcal F$ (Ordnung $k=1$).
Gamma-Matrizen erfüllen $\{\gamma^\mu,\gamma^\nu\}=2g^{\mu\nu}\mathbf 1$.

\emph{Übersetzer (äußere Symmetrie).}
\[
(\rho_{\mathrm{ext}}(a,\Lambda)\psi)(x')
:= S(\Lambda)\,\psi(\Lambda^{-1}(x'-a)),
\quad
S(\Lambda)^{-1}\gamma^\mu S(\Lambda)=\Lambda^\mu{}_\nu\gamma^\nu.
\]

\emph{Expliziter Nachweis der Natürlichkeit.}
Setze $\psi'=\rho_{\mathrm{ext}}(a,\Lambda)\psi$.
Wir berechnen $D_{x'}\psi'(x')$:
\[
\partial'_\mu \psi'(x')
= \partial'_\mu\big(S(\Lambda)\psi(\Lambda^{-1}(x'-a))\big)
= S(\Lambda)\,(\Lambda^{-1})^\nu{}_\mu\,\partial_\nu\psi(x),
\]
da $S(\Lambda)$ konstant (vom Raumzeitpunkt unabhängig) ist.
Dann
\[
i\gamma^\mu\partial'_\mu\psi'(x')
= i\gamma^\mu S(\Lambda)\,(\Lambda^{-1})^\nu{}_\mu\,\partial_\nu\psi(x)
= i\,S(\Lambda)\, S(\Lambda)^{-1}\gamma^\mu S(\Lambda)\,(\Lambda^{-1})^\nu{}_\mu\,\partial_\nu\psi(x).
\]
Mit der Kovarianz der Gamma-Matrizen
$S(\Lambda)^{-1}\gamma^\mu S(\Lambda)=\Lambda^\mu{}_\rho\gamma^\rho$
und $(\Lambda^{-1})^\nu{}_\mu \Lambda^\mu{}_\rho=\delta^\nu{}_\rho$ folgt
\[
i\gamma^\mu\partial'_\mu\psi'(x')
= i\,S(\Lambda)\,\gamma^\nu\,\partial_\nu\psi(x).
\]
Somit
\[
D_{x'}\psi'(x')
= i\gamma^\mu\partial'_\mu\psi'(x')-m\psi'(x')
= S(\Lambda)\,\big(i\gamma^\nu\partial_\nu\psi(x)-m\psi(x)\big)
= \rho_{\mathrm{ext}}(a,\Lambda)\big(D_x\psi\big)(x').
\]
In Funktornotation:
\[
D_{\Lambda U+a}\circ J^1(\rho_{\mathrm{ext}}(a,\Lambda))
= \rho_{\mathrm{ext}}(a,\Lambda)\circ D_U.
\]
\emph{Folge (Lösungserhalt):}
Erfüllt $\psi$ die Dirac-Gleichung auf $U$, dann erfüllt
$\psi'=\rho_{\mathrm{ext}}(a,\Lambda)\psi$ sie auf $\Lambda U + a$;
die Form des Gesetzes bleibt identisch.

\bigskip
\subsubsection*{C. Fazit in der Sprach der Unterfunktoren}

In beiden Beispielen ist das \emph{Gesetz} der \emph{Nullstellen-Unterfunktor}
\[
\mathcal E \subset \mathfrak F,\qquad \mathcal E(U)=\{\Phi\in\mathfrak F(U)\mid D_U(j^k\Phi)=0\}.
\]
Die obigen Rechnungen zeigen die \emph{Natürlichkeit} von $D$:
\[
D_{\Lambda U+a}\circ J^k(\rho_{\mathrm{ext}}(a,\Lambda))
=\rho_{\mathrm{ext}}(a,\Lambda)\circ D_U,
\]
und damit die \emph{Kovarianz} des Unterfunktors:
\[
\rho_{\mathrm{ext}}(a,\Lambda)\big(\mathcal E(U)\big)=\mathcal E(\Lambda U+a).
\]
Genau dies ist die präzise (jettheoretische) Aussage:
\[
\boxed{\text{„Gesetz ist forminvariant unter Beschreibungswechseln“}}
\]
— als Kommutativität eines natürlichen Differentialoperators mit den
durch die Symmetriegruppe induzierten Übersetzern.

\pagebreak

\subsection{Explizites Beispiel mit konkreter Lösung: Forminvarianz unter Lorentz-Transformationen}

\emph{Ziel.}  
Wir betrachten zwei explizite Fälle — das Klein–Gordon-Feld und das Dirac-Feld — 
und zeigen durch konkrete Lösungen, dass die Gleichung nach einem 
Beschreibungswechsel (Translation + Lorentz-Transformation) \emph{forminvariant} bleibt, 
also dieselbe Gestalt behält und der transformierte Ausdruck wieder eine Lösung ist.

\bigskip
\subsubsection*{A. Klein–Gordon-Feld}

\textbf{Gesetz:}
\[
(\Box + m^2)\,\phi = 0, 
\qquad 
\Box = g^{\mu\nu}\partial_\mu\partial_\nu,
\qquad
g = \mathrm{diag}(1,-1,-1,-1).
\]

\textbf{Feldbündel:}
\(\mathcal F = M\times\C,\)
\(\mathfrak F(U)=C^\infty(U,\C).\)

\textbf{Operator:}
\[
D = \Box + m^2 : J^2\mathcal F \to \mathcal F.
\]

\textbf{Konkrete Lösung:}
\[
\phi(x) = e^{-ip\cdot x} = e^{-i(Et-\mathbf{p}\cdot\mathbf{x})},
\quad
E=\sqrt{|\mathbf p|^2+m^2},
\]
und
\[
(\Box + m^2)\phi = (-p_\mu p^\mu + m^2)\phi = 0.
\]

\textbf{Beschreibungswechsel:}  
Für $(a,\Lambda)\in \R^{1,3}\rtimes SO^+(1,3)$ setze
\[
x'=\Lambda x + a,
\qquad
\phi'(x') := \phi(\Lambda^{-1}(x'-a)).
\]
Einsetzen ergibt
\[
\phi'(x') 
= e^{-ip\cdot\Lambda^{-1}(x'-a)}
= e^{+i(\Lambda p)\cdot a}\, e^{-i(\Lambda p)\cdot x'}.
\]
Bis auf die irrelevante globale Phase $e^{+i(\Lambda p)\cdot a}$ ist
\(\phi'\) wieder eine ebene Welle mit neuem Impuls
\[
p' = \Lambda p.
\]

\textbf{Überprüfung der Forminvarianz:}
\[
(\Box' + m^2)\phi'
= \big(-(p')^2 + m^2\big)\phi' 
= ( -p_\mu p^\mu + m^2)\phi' 
= 0.
\]
Damit gilt
\[
D_{\Lambda U + a}\big(\rho_{\mathrm{ext}}(a,\Lambda)\phi\big)
= \rho_{\mathrm{ext}}(a,\Lambda)\big(D_U\phi\big),
\]
also:
\[
D\circ J^2(\rho_{\mathrm{ext}}(a,\Lambda)) = \rho_{\mathrm{ext}}(a,\Lambda)\circ D.
\]
Dies ist die explizite Realisierung der \emph{Natürlichkeit} und Forminvarianz.

\bigskip
\textbf{Spezifisches Beispiel: Boost entlang der $z$-Achse.}

\[
\Lambda(v) =
\begin{pmatrix}
\gamma & 0 & 0 & -\gamma v\\
0 & 1 & 0 & 0\\
0 & 0 & 1 & 0\\
-\gamma v & 0 & 0 & \gamma
\end{pmatrix},
\qquad
\gamma = (1-v^2)^{-1/2}.
\]
Für
\[
p=(E,0,0,k),
\quad
p'=\Lambda(v)p=(\gamma(E-vk),0,0,\gamma(k-vE)),
\]
gilt
\[
\phi'(t',z') = e^{-ip'\cdot x'} = e^{-i(\gamma(E-vk)t' - \gamma(k-vE)z')},
\]
und $(\Box' + m^2)\phi' = 0$.  
Die Gleichung bleibt formgleich.



\textbf{Gegeben:}
\[
\phi'(x') = \phi(\Lambda^{-1}(x' - a)),
\qquad
x' = \Lambda x + a,
\]
mit Lorentztransformation $\Lambda\in SO^+(1,3)$ und Translation $a\in \R^{1,3}$.
Wir wollen berechnen, wie der D’Alembert-Operator auf $\phi'$ wirkt.

\bigskip
\textbf{1. Erste Ableitung.}

Nach der Kettenregel:
\[
\partial'_\mu \phi'(x') 
= \frac{\partial x^\nu}{\partial x'^\mu}\, \partial_\nu \phi(x)
= (\Lambda^{-1})^\nu{}_\mu \, \partial_\nu \phi(x),
\]
wobei $x = \Lambda^{-1}(x' - a)$.

\textbf{2. Zweite Ableitung.}

Wende erneut $\partial'_\nu$ an:
\[
\partial'_\nu \partial'_\mu \phi'(x')
= (\Lambda^{-1})^\alpha{}_\nu (\Lambda^{-1})^\beta{}_\mu \,
\partial_\alpha \partial_\beta \phi(x),
\]
da $(\Lambda^{-1})^\alpha{}_\nu$ konstant ist (Lorentz-Transformation ist linear).

\textbf{3. Wirkung des D’Alembert-Operators.}

\[
\Box' \phi'(x') 
:= g^{\mu\nu} \partial'_\mu \partial'_\nu \phi'(x')
= g^{\mu\nu} (\Lambda^{-1})^\alpha{}_\mu (\Lambda^{-1})^\beta{}_\nu \,
\partial_\alpha \partial_\beta \phi(x).
\]

\textbf{4. Lorentz-Invarianz der Metrik.}

Für $\Lambda\in SO(1,3)$ gilt
\[
g^{\mu\nu} (\Lambda^{-1})^\alpha{}_\mu (\Lambda^{-1})^\beta{}_\nu
= g^{\alpha\beta},
\]
denn per Definition: \(\Lambda^T g \Lambda = g \iff g = \Lambda^{-T} g \Lambda^{-1}\).

\textbf{Damit folgt:}
\[
\Box'\phi'(x') = g^{\alpha\beta} \partial_\alpha \partial_\beta \phi(x)
= (\Box \phi)(x).
\]

\textbf{5. Zusammen mit der Masse.}

Der Operator $(\Box + m^2)$ wirkt also kommutativ mit der Feldtransformation:
\[
(\Box' + m^2)\phi'(x') 
= (\Box + m^2)\phi(x).
\]
Daraus folgt sofort:

\[
(\Box + m^2)\phi = 0
\quad\Longrightarrow\quad
(\Box' + m^2)\phi' = 0.
\]

\textbf{6. Kontrolle am konkreten Beispiel.}

Setze $\phi(x)=e^{-ip\cdot x}$, dann:
\[
\phi'(x')=\phi(\Lambda^{-1}(x'-a)) = e^{-ip\cdot\Lambda^{-1}(x'-a)}
= e^{+i(\Lambda p)\cdot a}\, e^{-i(\Lambda p)\cdot x'}.
\]
Daher:
\[
\partial'_\mu \phi'(x') = -i (\Lambda p)_\mu\, \phi'(x'),
\qquad
\partial'_\mu \partial'_\nu \phi'(x') 
= -(\Lambda p)_\mu(\Lambda p)_\nu\, \phi'(x').
\]
Kontraktion:
\[
\Box'\phi'(x') = - g^{\mu\nu}(\Lambda p)_\mu(\Lambda p)_\nu\, \phi'(x')
= -p_\mu p^\mu\, \phi'(x'),
\]
da \(p'_\mu = \Lambda_\mu{}^\nu p_\nu\) und \(p'^2=p^2\).
Somit:
\[
(\Box'+m^2)\phi' = (-p'^2+m^2)\phi' = 0.
\]

\bigskip
\textbf{Fazit:}
\[
\boxed{
(\Box + m^2)\phi = 0
\quad\Longrightarrow\quad
(\Box' + m^2)\phi' = 0,
\qquad
\phi'(x')=\phi(\Lambda^{-1}(x'-a)).
}
\]
Damit ist die \emph{Forminvarianz} des Klein–Gordon-Gesetzes unter 
Translations- und Lorentz-Transformationen explizit bestätigt.


\bigskip
\subsubsection*{B. Dirac-Feld}

\textbf{Gesetz:}
\[
(i\gamma^\mu\partial_\mu - m)\,\psi = 0,
\qquad
\{\gamma^\mu,\gamma^\nu\} = 2g^{\mu\nu}\mathbf 1.
\]

\textbf{Feldbündel:}
\(\mathcal F = M\times\C^4,\)
\(\mathfrak F(U)=C^\infty(U,\C^4).\)

\textbf{Operator:}
\[
D = i\gamma^\mu\partial_\mu - m : J^1\mathcal F \to \mathcal F.
\]

\textbf{Konkrete Lösung:}
\[
\psi(x) = u_s(\mathbf p)\,e^{-ip\cdot x},
\qquad
(\slashed{p}-m)u_s(\mathbf p)=0,
\]
wobei
\[
u_s(\mathbf p)
=\sqrt{\frac{E+m}{2m}}
\begin{pmatrix}
\chi_s\\[2pt]
\frac{\boldsymbol{\sigma}\cdot\mathbf p}{E+m}\chi_s
\end{pmatrix},
\qquad
\chi_\uparrow=\binom10,\ \chi_\downarrow=\binom01.
\]

\textbf{Beschreibungswechsel:}  
Für $(a,\Lambda)$ definiere den Übersetzer
\[
\psi'(x') = S(\Lambda)\psi(\Lambda^{-1}(x'-a)),
\quad
S(\Lambda)^{-1}\gamma^\mu S(\Lambda)=\Lambda^\mu{}_\nu\gamma^\nu.
\]
Einsetzen der Lösung ergibt
\[
\psi'(x') = S(\Lambda)u_s(\mathbf p)\,e^{-ip\cdot\Lambda^{-1}(x'-a)}
= S(\Lambda)u_s(\mathbf p)\,e^{-i(\Lambda p)\cdot x'}\,e^{+i(\Lambda p)\cdot a}.
\]
Definiere
\[
p'=\Lambda p,
\qquad
u'_s(\mathbf p'):=S(\Lambda)u_s(\mathbf p).
\]
Dann ist
\[
\psi'(x')=u'_s(\mathbf p')\,e^{-ip'\cdot x'},
\]
wobei $(\slashed{p'}-m)u'_s(\mathbf p')=0$.

\textbf{Überprüfung der Forminvarianz:}
\[
(i\gamma^\mu\partial'_\mu - m)\psi'
= i\gamma^\mu S(\Lambda)(\Lambda^{-1})^\nu{}_\mu \partial_\nu \psi(\Lambda^{-1}x')
- mS(\Lambda)\psi(\Lambda^{-1}x').
\]
Mit $\gamma^\mu (\Lambda^{-1})^\nu{}_\mu = S(\Lambda)\gamma^\nu S(\Lambda)^{-1}$ folgt
\[
(i\gamma^\mu\partial'_\mu - m)\psi'
= S(\Lambda)(i\gamma^\nu\partial_\nu - m)\psi(\Lambda^{-1}x')
= S(\Lambda)\big(D\psi\big)(\Lambda^{-1}x')
= 0.
\]
Also:
\[
D_{\Lambda U + a}\big(\rho_{\mathrm{ext}}(a,\Lambda)\psi\big)
= \rho_{\mathrm{ext}}(a,\Lambda)\big(D_U\psi\big),
\]
d.h. $D$ ist natürlich und die Dirac-Gleichung formgleich.

\bigskip
\textbf{Explizit: Ruhespinor + Boost entlang der $z$-Achse.}

Für $p=(m,0,0,0)$ und $s=\uparrow$:
\[
u_\uparrow(\mathbf 0)=
\begin{pmatrix}1\\0\\1\\0\end{pmatrix},
\qquad
\psi(x)=u_\uparrow(\mathbf 0)\,e^{-imt}.
\]
Boost mit Geschwindigkeit $v$ in $z$-Richtung:
\[
S(\Lambda)=\exp\!\Big(-\frac{\eta}{2}\gamma^0\gamma^3\Big),
\quad
\tanh\eta=v,\quad
\gamma=\cosh\eta.
\]
Dann:
\[
u'_\uparrow(\mathbf p')
=S(\Lambda)u_\uparrow(\mathbf 0)
=\sqrt{\frac{E+m}{2m}}
\begin{pmatrix}
1\\[2pt]0\\[2pt]\frac{p'_z}{E+m}\\[2pt]0
\end{pmatrix},
\quad
p'=(\gamma m,0,0,\gamma v m).
\]
Die transformierte Lösung:
\[
\psi'(t',z')=u'_\uparrow(\mathbf p')\,e^{-i(\gamma m\, t' - \gamma v m\, z')},
\]
erfüllt wieder
\[
(i\gamma^\mu\partial'_\mu - m)\psi'=0.
\]

\bigskip
\textbf{Fazit:}
\begin{itemize}
\item Die Übersetzungsabbildung $(a,\Lambda)$ transportiert die Lösung auf eine neue Lösung.
\item Der Operator $D$ kommutiert mit dieser Übersetzung auf Jet-Ebene:
\[
D_{\Lambda U + a}\circ J^k(\rho_{\mathrm{ext}}(a,\Lambda))
= \rho_{\mathrm{ext}}(a,\Lambda)\circ D_U.
\]
\item Die \emph{Form} der Gleichung bleibt identisch; nur die Parameter 
(Frequenz, Impuls, Spinor) transformieren.
\end{itemize}
Dies ist die konkrete, berechnete Veranschaulichung des Prinzips:
\[
\boxed{
\text{„Natürlichkeit eines Differentialoperators“}
\;\Longleftrightarrow\;
\text{Forminvarianz des physikalischen Gesetzes.}
}
\]

\pagebreak

\subsection{Abstrakte Formulierung physikalischer Modelle}

\emph{1. Vorbereitende Kategorien.}

Sei $\mathbf{Man}$ die Kategorie der glatten Mannigfaltigkeiten und glatten Abbildungen.  
Sei $\mathbf{Bund}$ die Kategorie der glatten Vektorbündel (oder allgemeiner: Faserbündel) und glatten Bündelmorphismen.  

Wir betrachten einen allgemeinen Funktor
\[
\mathcal{F}: \mathbf{Man}^{\mathrm{op}} \longrightarrow \mathbf{Bund},
\]
der jeder Mannigfaltigkeit $M$ den „Raum der Felder“ $\mathcal{F}(M)\to M$ zuordnet.  
Abschnitte $\phi\in \Gamma(\mathcal{F}(M))$ sind potentielle \emph{Beschreibungen} des Systems auf der Arena $M$.

\bigskip
\emph{2. Strukturelle Daten.}

Ein \emph{physikalisches Modell} ist ein Tupel
\[
\mathsf{Sys}
= (\mathcal{A}, G_{\mathrm{ext}}, G_{\mathrm{int}},
  \mathcal{F}, \rho_{\mathrm{ext}}, \rho_{\mathrm{int}},
  \mathcal{E}, \mathcal{O})
\]
mit folgenden abstrakten Komponenten:

\begin{enumerate}
\item $\mathcal{A}$: ein Objekt einer geeigneten geometrischen Kategorie $\mathbf{Geom}$ (z.\,B.\ Mannigfaltigkeit mit Struktur).  
Dies ist die \emph{kinematische Arena}.

\item $G_{\mathrm{ext}}$: eine Lie-Gruppe mit einer treuen Aktion
\[
\alpha_{\mathrm{ext}}: G_{\mathrm{ext}} \curvearrowright \mathcal{A}
\]
in der Kategorie $\mathbf{Geom}$.  
Diese Aktion beschreibt die zulässigen Endomorphismen der Arena.

\item $G_{\mathrm{int}}$: eine Lie-Gruppe, deren Elemente als interne Automorphismen der Faserstruktur wirken; formal
\[
\alpha_{\mathrm{int}}: G_{\mathrm{int}} \curvearrowright \mathcal{F},
\]
wobei $\pi\circ \alpha_{\mathrm{int}}(g)=\pi$.

\item $\mathcal{F}$: ein Objekt in $\mathbf{Bund}$, ein glattes Bündel $\pi:\mathcal{F}\to\mathcal{A}$,
dessen Fasern $F_x$ den lokalen Zustandsraum bilden.

\item $\rho_{\mathrm{ext}}$ und $\rho_{\mathrm{int}}$: Gruppenhomomorphismen
\[
\rho_{\mathrm{ext}}: G_{\mathrm{ext}} \longrightarrow \Aut_{\mathbf{Bund}}(\mathcal{F}),
\qquad
\rho_{\mathrm{int}}: G_{\mathrm{int}} \longrightarrow \Aut_{\mathbf{Bund}}(\mathcal{F}),
\]
die die Wirkungen auf dem Bündel implementieren.  
Beide Wirkungen kommutieren in geeigneter Weise (semi-direktes Produkt oder Produktstruktur).

\item $\mathcal{E}\subset \Gamma(\mathcal{F})$: eine Untervarietät oder ein Unterfunktor, der die
zulässigen Feldkonfigurationen bestimmt (z.\,B.\ Lösungsraum einer Differentialrelation).  
Formal: $\mathcal{E}$ ist ein Unterfunktor von $\Gamma\circ\mathcal{F}$,
\[
\mathcal{E}(U) \subset \Gamma(U,\mathcal{F}|_U)
\quad\text{für jedes }U\subseteq\mathcal{A}.
\]

\item $\mathcal{O}$: eine strukturierte Kategorie oder Algebra von Funktionalen auf $\mathcal{E}$,
\[
\mathcal{O} = \mathrm{Hom}_{\mathbf{Set}}(\mathcal{E}, \mathbb{K})
\]
mit geeigneter algebraischer oder topologischer Struktur (*-Algebra, symplektischer Raum etc.).
\end{enumerate}

\bigskip
\emph{3. Beschreibungen als lokale Funktorobjekte.}

Für $U\subseteq\mathcal{A}$ offen sei
\[
\mathsf{Desc}(U) = \Gamma(U,\mathcal{F}|_U),
\]
die Kategorie (bzw.\ Menge) lokaler Feldkonfigurationen auf $U$.  
Die Einschränkungsabbildungen $r_{U,V}: \mathsf{Desc}(V)\to \mathsf{Desc}(U)$
für $U\subset V$ machen $U\mapsto\mathsf{Desc}(U)$ zu einem Präschemen (präfaserbündelartigen Funktor).

\bigskip
\emph{4. Transformationen.}

Definiere eine Kategorie der \emph{zulässigen Transformationen}:
\[
\mathbf{Transf} = G_{\mathrm{ext}}\times G_{\mathrm{int}},
\]
deren Objekte die Einheitsgruppe ist und deren Morphismen die Paare $(\Lambda,\gamma)$ mit
Kombinationsgesetz $(\Lambda_1,\gamma_1)\circ(\Lambda_2,\gamma_2)=(\Lambda_1\Lambda_2,\gamma_1\gamma_2)$.
Sie wirken auf $\mathsf{Desc}(\mathsf{Sys})$ als
\[
(\Lambda,\gamma)\cdot \phi
= \rho_{\mathrm{ext}}(\Lambda)\,\rho_{\mathrm{int}}(\gamma)\,\phi.
\]

\bigskip
\emph{5. Gruppoid der Beschreibungen.}

Die Daten
\[
\mathsf{Desc}(\mathsf{Sys})
=(\mathrm{Obj}=\{\text{lokale Felder}\},
 \mathrm{Mor}=\mathbf{Transf})
\]
bilden ein Gruppoid.  
Jedes Objekt $\phi$ besitzt als Isomorphiegruppe die Stabilitätsuntergruppe
\[
\mathrm{Stab}(\phi)
=\{(\Lambda,\gamma)\mid (\Lambda,\gamma)\cdot\phi=\phi\}.
\]
Der Orbitraum $\mathcal{E}/\mathbf{Transf}$ ist die Menge der \emph{physikalischen Zustände}.

\bigskip
\emph{6. Kovarianz und Invarianz.}

Eine Gleichungsmenge $\mathcal{E}\subset \Gamma(\mathcal{F})$ heißt
\emph{kovariant}, falls
\[
(\Lambda,\gamma)\cdot \mathcal{E} = \mathcal{E}
\quad\forall (\Lambda,\gamma)\in\mathbf{Transf}.
\]
Eine Observable $O\in\mathcal{O}$ heißt \emph{invariant}, falls
\[
O(\phi)=O\!\big((\Lambda,\gamma)\cdot\phi\big)
\quad\forall \phi\in\mathcal{E},\ (\Lambda,\gamma)\in\mathbf{Transf}.
\]
Diese Definitionen sind rein kategorial und unabhängig von der
konkreten Natur der Objekte.

\bigskip
\emph{7. Struktur der Funktorialität.}

Die Konstruktion $U\mapsto\mathsf{Desc}(U)$ und $(\Lambda,\gamma)\mapsto(\Lambda,\gamma)\cdot-$
definiert einen Funktor
\[
\mathsf{Desc}(\mathsf{Sys}): \mathbf{Transf}^{\mathrm{op}} \longrightarrow \mathbf{Set},
\]
der jedem Transformationspaar den Automorphismus seiner Wirkung auf dem
Beschreibungssatz zuordnet.  Die Kovarianz der Gleichungen entspricht der
\emph{Natürlichkeit} des Unterfunktors $\mathcal{E}$ bzgl.\ dieser Wirkung:
\[
\mathcal{E}\circ(\Lambda,\gamma)
= \mathcal{E}.
\]

\bigskip
\emph{8. Funktor der Lösungen und Observablen.}

Der \emph{Lösungsfunktor}
\[
\mathcal{S}ol: \mathbf{Transf}^{\mathrm{op}} \longrightarrow \mathbf{Set},
\qquad
(\Lambda,\gamma)\longmapsto
\{\phi\in\mathcal{E}\mid (\Lambda,\gamma)\cdot\phi\in\mathcal{E}\}
\]
ist ein $G_{\mathrm{ext}}\times G_{\mathrm{int}}$-äquivarianter Funktor.  

Observablen sind dann \emph{natürliche Transformationen}
\[
O:\mathcal{S}ol \Rightarrow \underline{\mathbb{K}},
\]
wobei $\underline{\mathbb{K}}$ der konstante Funktor auf dem Skalarkörper ist.
Invarianz entspricht der Natürlichkeitseigenschaft
\[
O_{\mathcal{E}} = O_{\mathcal{E}}\circ(\Lambda,\gamma).
\]

\bigskip
\emph{9. Abstrakte Lesart.}

Das formale Gerüst liefert folgende begriffliche Ebenen:
\begin{itemize}
\item Die \emph{Arena} $\mathcal{A}$: Träger der kinematischen Struktur.
\item Die \emph{Gruppen} $G_{\mathrm{ext}}, G_{\mathrm{int}}$: Generatoren zulässiger Transformationen.
\item Das \emph{Feldbündel} $\mathcal{F}$: Objekt, auf dem diese Gruppen wirken.
\item Die \emph{Darstellungen} $\rho_{\mathrm{ext}},\rho_{\mathrm{int}}$: Vermittlung zwischen Struktur und Wirkung.
\item Die \emph{Gleichungen} $\mathcal{E}$: Selektionsprinzip innerhalb des Feldraums.
\item Die \emph{Observablen} $\mathcal{O}$: Funktoren auf den Lösungsfunktor, invariant unter Transformationen.
\item Die \emph{Invarianz- und Kovarianzprinzipien}: Natürlichkeitsbedingungen der beteiligten Funktoren.
\end{itemize}

\emph{10. Zusammenfassung.}

Ein physikalisches System ist damit ein \emph{objektives Strukturensemble}, bestehend aus
\begin{enumerate}
\item einem \emph{Bezugssystem-Raum} (Arena $\mathcal{A}$ und externe Symmetrie $G_{\mathrm{ext}}$),
\item einem \emph{internen Zustandsraum} (Faserstruktur und interne Symmetrie $G_{\mathrm{int}}$),
\item einem \emph{Feldfunktor} $\mathcal{F}$, auf dem beide Gruppen durch Darstellungen wirken,
\item einem \emph{Lösungsfunktor} $\mathcal{E}$, kovariant unter diesen Wirkungen,
\item und einer \emph{Observablen-Kategorie} $\mathcal{O}$, deren Objekte natürliche Transformationen auf $\mathcal{E}$ sind.
\end{enumerate}


\subsection{Abstrakter Modellrahmen: formelle Ausarbeitung}

\emph{Notation.}
Es seien $\mathbf{Man}$ die Kategorie glatter, parakompakter Mannigfaltigkeiten mit glatten Abbildungen
und $\mathbf{Bund}$ die $\,\mathbf{Man}$–faserte Kategorie glatter (reeller/komplexer) Vektorbündel
mit Bündelmorphismen über der Identität. 
Für $M\in\mathrm{Obj}(\mathbf{Man})$ bezeichne $\Gamma(E\to M)$ die glatten Abschnitte.
Für eine Lie-Gruppe $G$ bezeichne $BG$ die Ein-Objekt-Kategorie (Klassenifying-Stack-Sicht).

\bigskip
\emph{1. Arenen und äußere Symmetrien.}
\begin{itemize}
\item Eine \emph{kinematische Arena} ist ein Objekt $\mathcal A$ einer geometrischen Unterkategorie
$\mathbf{Geom}\subset\mathbf{Man}$ (z.\,B.\ orientierte, pseudo-Riemannsche Mannigfaltigkeiten).
\item Eine \emph{äußere Symmetrie} ist eine glatte Linksaktion einer Lie-Gruppe $G_{\mathrm{ext}}$ auf $\mathcal A$ in $\mathbf{Geom}$:
\[
\alpha_{\mathrm{ext}}:G_{\mathrm{ext}}\times \mathcal A \to \mathcal A,\qquad
(\Lambda,x)\mapsto \Lambda\cdot x,
\]
äquivalent zu einem Funktor $G_{\mathrm{ext}}\to\mathrm{Aut}_{\mathbf{Geom}}(\mathcal A)$.
\item Die Kategorie der \emph{Hintergrund-Morphismen} ist die von $G_{\mathrm{ext}}$ erzeugte Unterkategorie
$\mathbf{Back}(\mathcal A)$ mit Objekten $\mathcal A$ und Morphismen $\{\Lambda:\mathcal A\to \mathcal A\}$.
\end{itemize}

\bigskip
\emph{2. Innere Symmetrien und Bündelstruktur.}
\begin{itemize}
\item Eine \emph{innere Symmetriegruppe} ist eine Lie-Gruppe $G_{\mathrm{int}}$.
\item Ein \emph{Hauptbündel} $P\to \mathcal A$ mit Strukturgruppe $G_{\mathrm{int}}$ ist gegeben durch eine freie, eigentliche Rechtsaktion
$P\times G_{\mathrm{int}}\to P$ mit Quotient $\mathcal A$.
\item Für eine Darstellung $\rho_{\mathrm{int}}:G_{\mathrm{int}}\to GL(F_0)$ ist das \emph{assoziierte Feldbündel}
\[
\pi:\ \mathcal F := P\times_{\rho_{\mathrm{int}}} F_0 \longrightarrow \mathcal A
\]
definiert. Eine \emph{Feldkonfiguration} ist $\phi\in\Gamma(\mathcal F)$.
\item Die \emph{Eichgruppe auf $U\subset\mathcal A$} ist die Gruppe der $G_{\mathrm{int}}$–äquivarianten Automorphismen von $P|_U$ über $\mathrm{id}_U$:
\[
\mathcal G_{\mathrm{int}}(U):=\mathrm{Aut}_{G_{\mathrm{int}}}(P|_U\to U).
\]
Sie bildet einen Prägarben-Funktor $U\mapsto\mathcal G_{\mathrm{int}}(U)$.
\end{itemize}

\bigskip
\emph{3. Feldfunktor, Beschreibungen und Restriktionen.}
\begin{itemize}
\item Der \emph{Feldfunktor} ist ein (prä-)Garben-Funktor
\[
\mathfrak F:\ \mathrm{Open}(\mathcal A)^{\mathrm{op}}\longrightarrow \mathbf{Set},\qquad
U\longmapsto \mathfrak F(U):=\Gamma(\mathcal F|_U).
\]
\item Für $U\subset V$ ist die Restriktion $r_{U,V}:\mathfrak F(V)\to \mathfrak F(U)$ natürlich.
\item Eine \emph{Beschreibung} ist ein Paar $(U,\phi)$ mit $U\subset \mathcal A$ offen, $\phi\in\mathfrak F(U)$.
\end{itemize}

\bigskip
\emph{4. Aktionen auf Feldern (äußere und innere).}
\begin{itemize}
\item Die äußere Symmetrie wirkt (kovariant) über einen Bündel-Automorphismus-Funktor
\[
\rho_{\mathrm{ext}}:\ G_{\mathrm{ext}}\longrightarrow \mathrm{Aut}_{\mathbf{Bund}}(\mathcal F),\qquad
\phi\ \mapsto\ \rho_{\mathrm{ext}}(\Lambda)\cdot \phi,
\]
kompatibel mit der Basiswirkung $\alpha_{\mathrm{ext}}(\Lambda)$ (Kartentransform).
\item Die innere Symmetrie wirkt vertikal
\[
\rho_{\mathrm{int}}:\ \mathcal G_{\mathrm{int}}(-)\longrightarrow \mathrm{Aut}_{\mathbf{Bund}}(\mathcal F|-),\qquad
\phi\ \mapsto\ \rho_{\mathrm{int}}(\gamma)\cdot \phi,
\]
faserweise über der Identität auf $U$.
\item Zusammen definiert $(\Lambda,\gamma)$ eine \emph{Transformation}
\[
(\Lambda,\gamma)\cdot (U,\phi)\ :=\ \big(\alpha_{\mathrm{ext}}(\Lambda)U,\ \rho_{\mathrm{ext}}(\Lambda)\rho_{\mathrm{int}}(\gamma)\,\phi\big).
\]
\end{itemize}

\bigskip
\emph{5. Gruppoid der Beschreibungen.}
\begin{itemize}
\item Das \emph{Transformationsobjekt} ist das semidirekte Produkt
$G_{\mathrm{tot}}:=G_{\mathrm{ext}}\ltimes \mathcal G_{\mathrm{int}}(-)$ (als Garbe in Gruppen).
\item Das \emph{Beschreibungsgrouppoid} ist
\[
\mathsf{Desc}(\mathsf{Sys})\ :=\
\Big(\ \mathrm{Obj}=\big\{(U,\phi)\big\},\ 
\mathrm{Mor}\big((U,\phi)\to(U',\phi')\big)=\big\{(\Lambda,\gamma)\mid (\Lambda,\gamma)\cdot(U,\phi)=(U',\phi')\big\}\ \Big).
\]
\item Der \emph{Orbit} eines $\phi\in\mathfrak F(U)$ ist $G_{\mathrm{tot}}(U)\cdot\phi$.
\end{itemize}

\bigskip
\emph{6. Natürliche Differentialoperatoren und Kovarianz.}
\begin{itemize}
\item Ein (lokaler) \emph{natürlicher Differentialoperator} $D$ Ordnung $\le k$ ist eine natürliche Transformation
\[
D:\ J^k\mathfrak F \Longrightarrow \mathfrak E,
\]
zwischen dem $k$-Jetfunktor $J^k\mathfrak F$ und einem Zielbündel-Funktor $\mathfrak E$,
so dass für jede Hintergrundmorphismen $\Lambda$ ein Kommutativitätsdiagramm gilt:
\[
\begin{array}{ccc}
J^k\mathfrak F(U) & \xrightarrow{\ \ D_U\ \ } & \mathfrak E(U)\\
\downarrow J^k(\rho_{\mathrm{ext}}(\Lambda)) & & \downarrow \rho_{\mathrm{ext}}(\Lambda) \\
J^k\mathfrak F(\Lambda U) & \xrightarrow{\ \ D_{\Lambda U}\ \ } & \mathfrak E(\Lambda U)
\end{array}
\]
und analog für $\rho_{\mathrm{int}}$.
\item Eine \emph{Gleichungsmenge} ist der Nullstellenfunktor eines natürlichen Operators:
\[
\mathcal E(U):=\{\ \phi\in\mathfrak F(U)\mid D_U(j^k\phi)=0\ \}.
\]
\item \emph{Kovarianz der Gleichungen} bedeutet Natürlichkeit:
\[
(\Lambda,\gamma)\cdot \mathcal E(U)=\mathcal E(\Lambda U),\quad
\text{d.\,h.}\quad
D_{\Lambda U}\circ J^k(\rho(\Lambda,\gamma))=\rho(\Lambda,\gamma)\circ D_U.
\]
\end{itemize}

\bigskip
\emph{7. Variationaler Formalismus (präzise).}
\begin{itemize}
\item Eine \emph{Lagrangedichte} ist eine natürliche horizontale $n$-Form
$\mathcal L\in\Omega^{n,0}(J^k\mathcal F)$ im Variationsbikomplex.
\item Der \emph{Euler–Lagrange-Operator} ist eine natürliche Transformation
$EL(\mathcal L):J^{2k}\mathfrak F\Rightarrow \mathfrak E^\vee$ der Form
\[
\mathcal E(U)=\{\phi\mid EL(\mathcal L)_U(j^{2k}\phi)=0\},
\]
und ist \emph{automatisch kovariant}, falls $\mathcal L$ natürlich ist.
\item \emph{Noether 1}: Eine einparametrige natürliche Symmetrie
$(\Lambda_t,\gamma_t)$ der Lagrangefunktionalität erzeugt (unter Regularität) einen horizontal geschlossenen Noether-Strom $J$, d.\,h.
$d_h J = EL(\mathcal L)\cdot \delta_\xi\phi$.
\end{itemize}

\bigskip
\emph{8. Observablen als natürliche Transformationen.}
\begin{itemize}
\item Der \emph{Lösungsfunktor} $\mathcal S ol:\mathrm{Open}(\mathcal A)^{\mathrm{op}}\to\mathbf{Set}$,
$U\mapsto \mathcal E(U)$, ist $G_{\mathrm{tot}}$–äquivariant.
\item Ein \emph{(klassisches) Observablen-Funktor} ist eine natürliche Transformation
\[
\mathcal O:\ \mathcal S ol \Longrightarrow \underline{\mathbb K}
\quad(\mathbb K=\R\ \text{oder}\ \C),
\]
die zusätzlich $G_{\mathrm{tot}}$–invariant ist (Natürlichkeit relativ zur Aktion).
\item Algebraische Struktur (Poisson, *-Algebra) wird als Monoidal-/*-Struktur in der Zielkategorie codiert und über Natürlichkeit bewahrt.
\end{itemize}

\bigskip
\emph{9. Quotienten, Moduli, Stack-Sicht.}
\begin{itemize}
\item Die \emph{Moduli der Lösungen bis Transformation} ist kein gutes Mengenkquotient; korrekt ist das \emph{Quotienten-Gruppoid} bzw.\ der \emph{Quotienten-Stack}
\[
\big[\ \mathcal S ol / G_{\mathrm{tot}}\ \big],
\]
dessen Objekte Lösungen, Morphismen Eich-/Bezugstransformationen sind.
\item Invariante Observablen sind Funktoren auf diesem Quotienten-Stack.
\end{itemize}

\bigskip
\emph{10. Zusammenfassende Axiome.}
Ein Datensatz
\[
\mathsf{Sys}=(\mathcal A, G_{\mathrm{ext}}, G_{\mathrm{int}}, \mathcal F,
\rho_{\mathrm{ext}},\rho_{\mathrm{int}}, \mathcal E, \mathcal O)
\]
heißt \emph{wohlgeformtes (klassisches) Modell}, falls:
\begin{enumerate}
\item (Feldgarbe) $U\mapsto \mathfrak F(U)=\Gamma(\mathcal F|_U)$ ist eine Prägarbe mit Sheaf-Eigenschaft (Gluing/Uniqueness).
\item (Natürliche Operatoren) Es existiert ein natürlicher Differentialoperator $D$ mit $\mathcal E=\ker D$.
\item (Kovarianz) Für alle $(\Lambda,\gamma)$ gilt Natürlichkeit
$D\circ J^k(\rho(\Lambda,\gamma))=\rho(\Lambda,\gamma)\circ D$.
\item (Gruppoid) $\mathsf{Desc}(\mathsf{Sys})$ ist ein wohldefiniertes Transformationsgruppoid, der Lösungsfunktor ist $G_{\mathrm{tot}}$–äquivariant.
\item (Observablen) $\mathcal O$ ist eine Kategorie natürlicher Transformationen 
$\mathcal S ol\Rightarrow\underline{\mathbb K}$, monoidal/*-stabil und $G_{\mathrm{tot}}$–invariant.
\item (Variational optional) Falls $\mathcal L$ existiert, ist $EL(\mathcal L)=0$ äquivalent zu $\mathcal E$, und Symmetrien induzieren Noether-Ströme.
\end{enumerate}

\bigskip
\emph{11. Diagrammatische Übersicht.}
\[
\begin{array}{ccccc}
G_{\mathrm{ext}}\times \mathcal G_{\mathrm{int}}(U) 
& \curvearrowright & \mathfrak F(U) & \xrightarrow{\ D_U\ } & \mathfrak E(U)\\[0.4em]
\downarrow & & \downarrow r_{U,V} & & \downarrow r_{U,V}\\[0.4em]
G_{\mathrm{ext}}\times \mathcal G_{\mathrm{int}}(V) 
& \curvearrowright & \mathfrak F(V) & \xrightarrow{\ D_V\ } & \mathfrak E(V)
\end{array}
\qquad
\begin{array}{c}
\mathcal E(U)=\ker D_U \\[0.3em]
\rotatebox{90}{$\subset$}\\[0.3em]
\mathfrak F(U)
\end{array}
\]
Natürlichkeit $\Leftrightarrow$ Kommutativität der Rechtecke; 
Kovarianz $\Leftrightarrow$ Stabilität von $\ker D$ unter der linken Wirkung.

\bigskip
\emph{12. Einsatzpunkte für weitere Präzisierungen.}
\begin{itemize}
\item (Jet-Bündel) Explizite Ordnung/Regularität von $D$ via $J^k\mathcal F$.
\item (Hyperbolizität) Wohlgestelltheit des Cauchy-Problems: symmetrisch-/normal-hyperbolische Klassen.
\item (Bikomplex) Variationsbikomplex, homologische Ableitungen, Horiz./Vert.-Differentiale.
\item (Quantisierung) Algebraisierung der $\mathcal O$ (Poisson $\to$ *-Algebra), lokalkovariante QFT als Funktor $\mathbf{Loc}\to\mathbf{Alg}$.
\end{itemize}


\subsection{Ein formaler Rahmen für Beschreibungen physikalischer Systeme}

\emph{Daten eines Modells.}
Ein (klassisch-feldtheoretisches) physikalisches System wird durch folgende Basisdaten beschrieben:
\[
\mathsf{Sys}
\;=\;
(\mathcal{A},\,G_{\mathrm{ext}},\,G_{\mathrm{int}},\,\mathcal{F},\,\rho_{\mathrm{ext}},\,\rho_{\mathrm{int}},\,\mathcal{E},\,\mathcal{O})
\]
\begin{itemize}
\item $\mathcal{A}$: \emph{kinematische Arena} (z.\,B. Raumzeit $(M,g)$; in SR: $M=\R^{1,3}$ mit Minkowski-Metrik $g$).
\item $G_{\mathrm{ext}}$: \emph{äußere Symmetriegruppe} der Arena (z.\,B. Poincaré-Gruppe, Diff$(M)$, Isometriegruppe).
\item $G_{\mathrm{int}}$: \emph{innere (Eich-)Symmetriegruppe} (z.\,B. $U(1),\, SU(2),\, SU(3)$).
\item $\mathcal{F}$: \emph{Feldbündel/Träger} der Freiheitsgrade (Abschnitte sind Feldkonfigurationen).
\item $\rho_{\mathrm{ext}},\rho_{\mathrm{int}}$: \emph{Darstellungen/Anbindungen} von $G_{\mathrm{ext}}$ und $G_{\mathrm{int}}$ auf $\mathcal{F}$ (z.\,B. Spin-Darstellung, Eichdarstellung).
\item $\mathcal{E}\subset \Gamma(\mathcal{F})$: \emph{Gleichungsmenge} (Lösungsmenge) als Untermenge/Untervarietät der Feldkonfigurationen (z.\,B. Euler–Lagrange-Gleichungen).
\item $\mathcal{O}$: \emph{Observablenraum} (z.\,B. Funktionale $\mathcal{E}\to \R/\C$, algebraische Struktur, *-Algebra).
\end{itemize}

\emph{Beschreibungen und Wechsel.}
\begin{itemize}
\item Eine \emph{Beschreibung} ist ein Paar $(U,\phi)$, wobei $U\subseteq \mathcal{A}$ (z.\,B. Kartenbereich/Trivialisierung) und
$\phi\in \Gamma(\mathcal{F}|_U)$ eine lokale Feldkonfiguration ist.
\item Ein \emph{Wechsel der Beschreibung} ist ein Morphismus
\[
(\Lambda,\gamma): (U,\phi)\longrightarrow (U',\phi')
\]
bestehend aus einer \emph{äußeren} Transformation $\Lambda\in G_{\mathrm{ext}}$ (Referenzrahmenwechsel) und einer \emph{inneren} Transformation $\gamma\in \mathcal{G}_{\mathrm{int}}(U)$
(lokale Eichtransformation), so dass
\[
\phi' \;=\; \rho_{\mathrm{ext}}(\Lambda)\;\rho_{\mathrm{int}}(\gamma)\;\phi
\quad\text{auf}\quad U\cap U'.
\]
\end{itemize}
Die Gesamtheit der Beschreibungen und Wechsel bildet eine \emph{Gruppoid}-Struktur
\[
\mathsf{Desc}(\mathsf{Sys})
\]
(Objekte: Beschreibungen, Pfeile: zulässige Wechsel). Physikalische \emph{Zustände} sind dann \emph{Orbits} dieser Gruppoidwirkung
(„die gleiche Physik, unterschiedliche Beschreibungen“).

\emph{Kovarianz und Invarianz.}
\begin{itemize}
\item \emph{Kovarianz} der Gleichungen bedeutet:
\[
(\Lambda,\gamma)\cdot \mathcal{E} \;=\; \mathcal{E}
\quad\text{für alle}\quad (\Lambda,\gamma).
\]
Also: Die Lösungsmenge ist \emph{stabil} unter Referenzrahmen- und Eichwechseln (Gleichungsform bleibt erhalten).
\item \emph{Invarianz} einer Observablen $O\in\mathcal{O}$ bedeutet:
\[
O(\phi) \;=\; O\!\big((\Lambda,\gamma)\cdot\phi\big)
\quad\text{für alle}\quad \phi\in \mathcal{E},\;(\Lambda,\gamma).
\]
Physikalisch messbare Größen sind \emph{gruppoid-invariant}.
\end{itemize}

\emph{Äußere vs. innere Symmetrien.}
\begin{itemize}
\item $G_{\mathrm{ext}}$ wirkt auf der Arena (Koordinaten/Frames) und auf Tensorspuren (z.\,B. $V=\R^{1,3}$, Poincaré-Wirkung).
\item $G_{\mathrm{int}}$ wirkt \emph{fasertreu} auf $\mathcal{F}$ (interne Freiheitsgrade), typischerweise als Strukurgruppe eines Hauptbündels $P\to \mathcal{A}$, mit $\mathcal{F}$ assoziiert: $\mathcal{F}=P\times_\rho F$.
\end{itemize}

\emph{Darstellungen (Rollenklärung).}
\begin{itemize}
\item $\rho_{\mathrm{ext}}: G_{\mathrm{ext}}\to \mathrm{Aut}(\mathcal{F})$ implementiert den Wechsel der Bezugssysteme (z.\,B. Spin-Lift von Lorentz auf Spinoren).
\item $\rho_{\mathrm{int}}: G_{\mathrm{int}}\to \mathrm{Aut}(\mathcal{F})$ implementiert interne Symmetrien (z.\,B. Phasen für $U(1)$).
\item \emph{Equivarianzbedingung}: Strukturobjekte (z.\,B. Dirac-$\gamma$) müssen die natürliche Kompatibilität erfüllen:
\[
\rho_{\mathrm{ext}}(\Lambda)^{-1}\,\gamma^\mu\,\rho_{\mathrm{ext}}(\Lambda)
\;=\; \Lambda^\mu{}_\nu \,\gamma^\nu,
\]
damit $i\gamma^\mu\partial_\mu$ als natürlicher (erster Ordnung) Operator zwischen Beschreibungen fungiert.
\end{itemize}

\emph{Lagrange-zu-Gleichungen.}
Gegeben eine dichte Lagrange-Form $\mathcal{L}(\phi,\partial\phi; \text{Hintergrund})$, definiere die Euler–Lagrange-Gleichungen
$\mathrm{EL}(\mathcal{L})=0$.
Kovarianz/Invarianz auf Lagrange-Ebene bedeutet Naturtreue:
\[
\mathcal{L}\big((\Lambda,\gamma)\cdot\phi\big)
\;=\; \mathcal{L}(\phi) \;+\; \partial_\mu(\cdots)^\mu
\]
(bis auf Totalableitungen). Kontinuierliche Symmetrien $\Rightarrow$ Noether-Ströme.

\emph{Kategorien-/Funktor-Sicht.}
\begin{itemize}
\item Kategorie der Beschreibungen $\mathsf{Desc}(\mathsf{Sys})$ und Funktor der \emph{Lösungen}
\[
\mathcal{S}ol:\;\mathsf{Desc}(\mathsf{Sys}) \longrightarrow \mathsf{Set},\qquad
(U,\phi)\longmapsto \{\psi\in\Gamma(\mathcal{F}|_U)\mid \psi\ \text{löst}\ \mathcal{E}\}.
\]
\item \emph{Kovarianz} $\Leftrightarrow$ \emph{Natürlichkeit} von $\mathcal{S}ol$ bzgl. Morphismen (Wechseln).
\item Observablen als natürl. Transformationen $\mathcal{S}ol\Rightarrow \underline{\C}$, die auf Isomorphismen (Wechsel) invariant sind.
\end{itemize}

\emph{Kommutatives Diagramm der Symmetrien.}
\[
\begin{array}{ccc}
G_{\mathrm{ext}} \times G_{\mathrm{int}}
 & \xrightarrow{\ \rho_{\mathrm{ext}}\times\rho_{\mathrm{int}} \ } &
\mathrm{Aut}(\mathcal{F}) \\[0.4em]
\downarrow &  & \downarrow \mathrm{induz.} \\[0.4em]
\mathrm{Aut}(\mathcal{A})
 & \xrightarrow{\ \text{Pull/Push}\ } &
\mathrm{Aut}(\Gamma(\mathcal{F}))
\end{array}
\]
Stabilität von $\mathcal{E}$ unter der rechten Wirkung ist die präzise Form von \emph{Gleichungs-Kovarianz}.

\bigskip

\subsubsection*{Instanziierung: Dirac-Feld im Minkowski-Raum (QFT-Setting)}

\emph{Arena und Gruppen.}
\[
\mathcal{A}=(M,g)=(\R^{1,3},\,\mathrm{diag}(1,-1,-1,-1)),\quad
G_{\mathrm{ext}}=\mathcal{P}=\R^{1,3}\rtimes SO^+(1,3),\quad
G_{\mathrm{int}}=U(1).
\]

\emph{Feldträger und Darstellungen.}
\[
\mathcal{F}= M\times \Sigma,\quad \Sigma=\C^4,\quad
\rho_{\mathrm{ext}}|_{SO^+(1,3)}=\sigma:\Spin(1,3)\to GL(\Sigma),
\]
\[
\text{Clifford: }\ \{\gamma^\mu,\gamma^\nu\}=2g^{\mu\nu}\mathbf{1},
\qquad
\rho_{\mathrm{int}}(e^{i\alpha})\psi = e^{-iq\alpha}\psi.
\]
Equivarianz der $\gamma$:
\[
\sigma(a)^{-1}\gamma^\mu\sigma(a)=\Lambda^\mu{}_\nu\gamma^\nu,\qquad \Lambda=\rho(a).
\]

\emph{Dynamik (frei bzw. minimal gekoppelt).}
\[
\mathcal{L}_{\mathrm{Dirac}}=\bar\psi\,(i\gamma^\mu\partial_\mu - m)\,\psi,
\qquad
\text{bzw.}\quad
\mathcal{L}_{\mathrm{QED}}=\bar\psi\,(i\gamma^\mu D_\mu - m)\psi - \tfrac{1}{4}F_{\mu\nu}F^{\mu\nu},
\]
\[
D_\mu=\partial_\mu + i q A_\mu,\quad
\psi \mapsto e^{-iq\alpha(x)}\psi,\quad
A_\mu\mapsto A_\mu+\partial_\mu\alpha.
\]
Euler–Lagrange:
\[
(i\gamma^\mu\partial_\mu - m)\psi=0
\quad\text{bzw.}\quad
(i\gamma^\mu D_\mu - m)\psi=0.
\]

\emph{Kovarianz/Invarianz.}
\begin{itemize}
\item \textbf{Äußere Symmetrie (Poincaré):} $x'=\Lambda x + a$, $\psi'(x')=S(\Lambda)\psi(x)$ mit $S=\sigma(a)\in \Spin(1,3)$,
\[
S^{-1}\gamma^\mu S=\Lambda^\mu{}_\nu\gamma^\nu,
\quad\Rightarrow\quad
(i\gamma^\mu\partial_\mu - m)\psi=0\ \Leftrightarrow\ (i\gamma^\mu\partial'_\mu - m)\psi'=0.
\]
\item \textbf{Innere Symmetrie ($U(1)$):}
$\psi\mapsto e^{-iq\alpha(x)}\psi$, $A\mapsto A+d\alpha$ lässt $(i\gamma^\mu D_\mu - m)\psi=0$ invariant.
\end{itemize}

\emph{Observablen.}
\[
j^\mu=\bar\psi\gamma^\mu\psi \quad (\text{Noether-Strom für }U(1)),\qquad
\partial_\mu j^\mu=0.
\]
Messbare Größen sind Funktionen auf dem Orbitraum $\mathcal{E}/(G_{\mathrm{ext}}\times \mathcal{G}_{\mathrm{int}})$.

\bigskip

\emph{Lesart in einem Satz.}
Bezugssysteme und Eichwahlen sind \emph{Beschreibungen} desselben physikalischen Zustands.
Die Menge aller Beschreibungen bildet ein Gruppoid, dessen Wirkung auf Feldern die
\emph{Gleichungsmenge} stabil lässt (Kovarianz). Physikalische Observablen sind Funktionen auf dem \emph{Orbitraum} (Invarianz). Darstellungen von $G_{\mathrm{ext}}$ (Spin-Lift) und $G_{\mathrm{int}}$ (Eichdarstellung) implementieren die Wechsel zwischen Beschreibungen; Strukturobjekte (z.\,B. $\gamma^\mu$) sind equivariant, sodass die Form der Gleichungen bewahrt bleibt.


\pagebreak
\printbibliography
\end{document}